\chapter{Singular backgrounds and vacuum data}
\label{v4-ch:singular-backgrounds}

The simplest singular current already separates a vacuum
from a scalar augmentation. At nonzero level \(\kappa\),
its first product satisfies
\[
                         J_{(1)}J=\kappa\mathbf1.
\]
Every first product in the scalar vertex algebra is zero.
A unital field map from this current algebra to the
scalars is therefore impossible. Nevertheless, the current
has a vacuum, a vacuum expectation and a geometric
realization with scalar empty coefficient. These are
different structures.

We first construct that realization from all collision
strata. We then allow infinitely many current species.
Completing each energy space retains the local field
operations, including quadratic fields with
Hilbert--Schmidt coefficients. A change of coordinate
can require self-contractions that are discontinuous in
this completion. This leads to two further coefficient
domains, one controlling projective tensor norms and
one retaining contraction values as independent central
coordinates.

Gaussian preparations give actual
normalized states in a specified Hilbert sector. Their
existence and their transport require their own estimates.
We also determine when a formal state parameter may be
identified with an insertion coordinate. The resulting
completed products retain convergent mode identities
on lower-bounded modules.

\section{A current with a vacuum and no scalar augmentation}
\label{v4-sec:singular-heisenberg}

Let \(k\) be a field of characteristic zero and fix
\(\kappa\in k^\times\). The state space of one current is
\[
 H_\kappa=k[x_1,x_2,\ldots],\qquad \mathbf1=1.
\]
Every state is a polynomial in finitely many variables.
Give \(x_m\) energy \(m\). For \(m>0\), define operators
\[
 a_{-m}f=x_mf,\qquad
 a_mf=\kappa m\frac{\partial f}{\partial x_m},
 \qquad a_0=0.
\]

They obey
\begin{equation}\label{v4-eq:singular-CCR}
 [a_m,a_n]=\kappa m\delta_{m+n,0},
 \qquad a_m\mathbf1=0\quad(m\geq0).
\end{equation}
The translation operator is
\[
 T=\sum_{m\geq1}m x_{m+1}\frac{\partial}{\partial x_m}.
\]
This sum is finite on each state. It satisfies
\(T\mathbf1=0\) and \([T,a_m]=-m a_{m-1}\).

Put
\[
 J(z)=\sum_{m\in\mathbb Z}a_mz^{-m-1},\qquad
 J^{[m]}(z)=\frac{\partial_z^{m-1}J(z)}{(m-1)!}
 \quad(m\geq1).
\]
Normal ordering, denoted by colons, moves negative modes
to the left of positive modes. Define the field of a
monomial by
\begin{equation}\label{v4-eq:singular-state-field}
 Y(x_{m_1}\cdots x_{m_r},z)
       =:J^{[m_1]}(z)\cdots J^{[m_r]}(z):,
 \qquad Y(\mathbf1,z)=\mathrm{id}.
\end{equation}
Linearity defines \(Y\) on \(H_\kappa\). In the
notation \(Y(a,z)b=\sum_n a_{(n)}b\,z^{-n-1}\),
the current state is \(J=x_1\).

\begin{proposition}\label{v4-prop:singular-heisenberg-fields}
These formulas make \(H_\kappa\) a local vertex algebra
with vacuum \(\mathbf1\) and translation \(T\).
For states of energies \(E,F\), the locality exponent
\(E+F\) suffices. The elementary contraction is
\begin{equation}\label{v4-eq:singular-derivative-contraction}
 J^{[m]}(z)J^{[n]}(w)\sim
 \frac{\kappa(-1)^{m-1}(m+n-1)!}
      {(m-1)!(n-1)!(z-w)^{m+n}}.
\end{equation}
Here \(\sim\) records only the negative powers of
\(z-w\).
\end{proposition}

\begin{proof}
Split \(J\) into its negative- and positive-mode parts.
Equation \eqref{v4-eq:singular-CCR} gives
\[
 [J_+(z),J_-(w)]
     =\kappa\sum_{m\geq1}m z^{-m-1}w^{m-1}
     =\iota_{z,w}\frac{\kappa}{(z-w)^2},
\]
where \(\iota_{z,w}\) means expansion in powers of
\(w/z\). Commuting the positive modes in one normally
ordered monomial past the negative modes in another
gives a sum over cross pairings. A pairing uses a factor
at most once. The unpaired factors remain normally
ordered. Differentiating the displayed contraction
gives \eqref{v4-eq:singular-derivative-contraction}.

This expansion is the Wick formula. A pair
of energies \(m,n\) contributes a pole of
order \(m+n\). Thus states of energies \(E,F\)
have total pole order at most \(E+F\).
The two orders of the fields expand the same rational
contractions in the opposite chambers. Multiplication
by \((z-w)^{E+F}\) removes their difference.

On the vacuum, the negative part of \(J\) is
\(\sum_{r\geq0}x_{r+1}z^r=e^{zT}x_1\).
Differentiation and the product rule give
\(Y(a,z)\mathbf1=e^{zT}a\) for every monomial.
The commutator relation with \(T\) gives
\([T,Y(a,z)]=\partial_zY(a,z)=Y(Ta,z)\).

For three monomials, every nested Wick expansion
enumerates the same pairings, with the same
uncontracted factors. After clearing the finite
pairwise pole bound, their coefficient identities
agree. The chamber expansions of this identity are
the vertex Jacobi identity. All its coefficients
on a fixed state use finitely many annihilators.
Thus lower truncation, the vacuum, translation and
locality identities hold on the stated polynomial
space.

These are also the explicit formulas of the
Heisenberg construction.\footnote{E.~Frenkel,
\emph{Vertex Algebras and Algebraic Curves},
\S\S2.1--2.4,
\href{https://arxiv.org/abs/math/0007054}{arXiv:math/0007054}.}
\end{proof}

For example,
\[
 Y(J,u)J=\kappa u^{-2}\mathbf1+(e^{uT}J)J,
 \qquad J_{(1)}J=\kappa\mathbf1,
 \qquad J_{(0)}J=0.
\]
The product on the right in the regular term is
polynomial multiplication of creation states.
It should not be substituted for every field product.

\begin{proposition}\label{v4-prop:singular-no-augmentation}
There is no unital vertex-algebra map \(H_\kappa\to k\),
where \(k\) has its constant fields and zero translation.
The polynomial constant-term functional
\(\epsilon_0:H_\kappa\to k\) still defines a normalized
vacuum expectation.
\end{proposition}

\begin{proof}
A field map would preserve the first product and the
vacuum. It would send \(J_{(1)}J\) to \(\kappa\),
but the first product of any two scalar states is
zero. This contradicts \(\kappa\ne0\).
The functional \(\epsilon_0\) is defined on every
polynomial and has \(\epsilon_0(\mathbf1)=1\).
Applying it to products of fields gives the Wick
vacuum coefficients. No preservation of every
field product is asserted for this functional.
\end{proof}

The construction extends to a finite-dimensional
space of species with a specified symmetric pairing.
For orthogonal level-one currents, contractions of
different species vanish. Their vacuum tensor product
is the tensor product of their vertex algebras. The
creation-state energy has Hilbert series
\[
                     \prod_{m\geq1}(1-q^m)^{-r}
\]
for \(r\) species. The coefficient of each power of
\(q\) is finite. An arithmetic energy weighting or
a chosen charge is additional data on this space.

\subsection{A finite-place coefficient field and its diagonal}

The same construction over \(\mathbb Q_p\)
has the nonzero formal pole
\(\kappa(z-w)^{-2}\).
Its insertion variable is a coordinate on
an auxiliary curve over \(\mathbb Q_p\).
It is independent of the \(p\)-adic
topology on its coefficient field.
An analytic interpretation requires specified
norms and convergence domains.

Two completions of a tensor product
illustrate why this distinction matters.
Set \(R=\mathbb Z_p\otimes_{\mathbb Z}\mathbb Z_p\),
with multiplication map \(\mu:R\to\mathbb Z_p\).
Define its \(p\)-adic completion by
\(\widehat R_p=\varprojlim_r R/p^rR\).

\begin{proposition}\label{v4-prop:singular-padic-tensor}
Multiplication induces
\(\widehat R_p\simeq\mathbb Z_p\).
The ordinary multiplication map
\(R\to\mathbb Z_p\) is not injective.
After inverting \(p\), its diagonal ideal
has nonzero conormal module.
\end{proposition}

\begin{proof}
For every \(r\geq1\), tensoring with
\(\mathbb Z/p^r\) gives
\[
 R/p^rR
   \cong(\mathbb Z_p/p^r)\otimes_{\mathbb Z}
                         (\mathbb Z_p/p^r)
   \cong\mathbb Z/p^r.
\]
These multiplication isomorphisms commute
with reduction, proving the completion formula.
No interchange of an ordinary tensor
with an inverse limit is used.

Choose \(a\in\mathbb Z_p\setminus\mathbb Q\).
In \(\mathbb Q_p\otimes_{\mathbb Q}\mathbb Q_p\),
the element \(a\otimes1-1\otimes a\) is nonzero.
Indeed a \(\mathbb Q\)-linear functional
\(\ell:\mathbb Q_p\to\mathbb Q\) with
\(\ell(1)=0,\ell(a)=1\) detects it.
Its antecedent in \(R\) is therefore nonzero,
while multiplication kills it.
The reduction formula also shows that
it vanishes in every \(R/p^rR\).

For the conormal assertion, choose \(a\)
transcendental over \(\mathbb Q\).
Such elements exist because \(\mathbb Z_p\)
is uncountable and algebraic numbers are countable.
Include \(a\) in a transcendence basis.
The derivation sending \(a\) to \(1\)
and all other basis elements to zero
extends across the separable algebraic extension
to \(\mathbb Q_p\).
Hence \(\Omega_{\mathbb Q_p/\mathbb Q}\ne0\).

If \(I\) is the kernel of multiplication
on \(\mathbb Q_p\otimes_{\mathbb Q}\mathbb Q_p\),
the defining universal derivation gives
\[
                         I/I^2\cong\Omega_{\mathbb Q_p/\mathbb Q}.
\]
It sends the class of \(1\otimes a-a\otimes1\)
to \(da\).
This proves nontriviality of the first
infinitesimal neighborhood of that diagonal.
\end{proof}

Thus completion in \(p\) and completion
in the diagonal ideal concern different directions.
On the auxiliary line \(\mathbb A^1_{\mathbb Q_p}\),
the diagonal instead has normal parameter
\(u=z-w\).
Its formal neighborhood has coefficient ring
\(\mathbb Q_p[z][[u]]\), and its principal
parts retain every \(u^{-r-1}du\).
The formal current pole uses this parameter.
The \(p\)-adic completion formula alone
cannot eliminate these operations.

\begin{proposition}\label{v4-prop:singular-common-coefficients}
Any family of characteristic-zero fields
\((k_v)\) has an abstract common field
extension over \(\mathbb Q\).
The formal Heisenberg algebra over
\(\mathbb Q\) also base changes to each \(k_v\).
\end{proposition}

\begin{proof}
Form the filtered union
\[
                  R_\infty=\varinjlim_{S\text{ finite}}
                                 \bigotimes_{v\in S,\mathbb Q}k_v,
\]
where the maps insert units.
Each finite tensor is a nonzero
\(\mathbb Q\)-vector space.
The transition maps are injective:
a linear functional on the new field
taking \(1\) to \(1\) gives a retraction.
Hence \(R_\infty\) is a nonzero commutative ring.

Choose a maximal ideal.
The quotient field receives a unital
map from every \(k_v\).
Each such map is injective, since its
source is a field.
The explicit mode formulas and their
identities commute with scalar extension,
proving the final assertion.
\end{proof}

This field construction makes no topological
or canonical choice.
For example, an abstract embedding
\(\mathbb Q_p\to\mathbb C\) cannot be continuous:
\(p^r\to0\) in \(\mathbb Q_p\), while
its images are the unbounded real numbers \(p^r\).
Likewise the adelic projection
\(\mathbb A_{\mathbb Q}\to\mathbb Q_p\)
has the opposite direction from scalar extension
of a \(\mathbb Q_p\)-algebra to the adeles.
A placewise module construction must specify
its actual rings, maps and restricted topology.
Hecke actions and automorphic comparison maps
are further structures on those coefficients.

\section{The residue coalgebra of the singular background}
\label{v4-sec:singular-residue-coalgebra}

Let \(X=\mathbb A^1_k\). We use the word carrier
\(\mathcal W_X\), its cut tensor, and the partition
embeddings \(i_\pi,j_\pi\) of
Chapter~\ref{v4-ch:regular-coefficient-models}.
Order the blocks of \(\pi\) by first occurrence.

For \(r=|\pi|\), put
\[
 M_\pi(H_\kappa)=\omega_{X^\pi}\otimes_k H_\kappa^{\otimes r}.
\]
Its right derivative in the coordinate of block \(C\)
is minus the scalar derivative minus \(T\) on its
\(H_\kappa\)-factor. All external coefficient tensors
in this construction are algebraic over \(k\).

Define
\begin{equation}\label{v4-eq:singular-residue-object}
 (C_\kappa)_I
   =\bigoplus_{\pi\in\Pi(I)}
      i_{\pi,*}j_{\pi,*}j_\pi^*M_\pi(H_\kappa)[|\pi|],
 \qquad (C_\kappa)_\varnothing=k.
\end{equation}
The state complex has zero differential. Its \(r\)-block
term lies in cohomological degree \(-r\).

For a merger of blocks \(i<j\), take
\(u=z_{C_i}-z_{C_j}\), retain \(z_{C_j}\), and place
the merged block in position \(i\). Replace its two
states by \(Y(v_i,u)v_j\), leaving the other block
states in their original relative order. The face is
\begin{equation}\label{v4-eq:singular-residue-face}
 D_{\pi\rho}(f\otimes\mathbf v)
   =(-1)^{j-2}\operatorname{pp}_u
      \bigl(f\otimes
       v_1\otimes\cdots\otimes Y(v_i,u)v_j
                    \otimes\cdots\widehat{v_j}\cdots\bigr).
\end{equation}

The principal part uses the full right normal module
\(e_m=u^{-m-1}du=\delta\partial_u^m/m!\).
Other coordinate differences are expanded normally
before taking the principal part. Normal transfer
extends the map to the diagonal direct image.
For each input, lower truncation and the finite pole
order of \(f\) make its output a finite normal sum.

\begin{theorem}\label{v4-thm:singular-residue-coalgebra}
The sum of the faces \eqref{v4-eq:singular-residue-face}
is a differential. The objects \eqref{v4-eq:singular-residue-object}
have coherent exceptional comparisons for monotone
collisions and form a connected coaugmented cut
coalgebra \(C_\kappa\).
\end{theorem}

\begin{proof}
Translation covariance of the field identifies the
input and output tangential connections. In the normal
direction the identity
\(e_m\partial_u=(m+1)e_{m+1}\), together with
\(Y(Ta,u)b=\partial_uY(a,u)b\), gives differential
linearity of every face.

Two consecutive faces either merge disjoint pairs or
merge three blocks. In the first case the coefficient
operations commute, and the two normal orientations
have opposite signs. In the second case expand all
three inputs by the Wick formula before taking any
principal parts. There is one finite collection of
contraction patterns. Each pattern is a simultaneous
coefficient with a finite product of pairwise
denominators and a formal regular numerator.

Locality permits its expansion in each of the three
collision orders. The oriented contacts therefore
obey \(L-R+N=0\), by
Lemma~\ref{v4-lem:regular-three-channel}, with the
uncontracted states as coefficients.
The normal-coordinate transfer in that lemma is
part of the comparison. These two cases exhaust the
square of the differential.

Let \(\alpha:I\twoheadrightarrow J\) be monotone.
Exceptional pullback kills a partition term unless
each \(\alpha\)-fiber lies in one of its blocks.
Every other term meets a diagonal removed by \(j_\pi\).
Every surviving partition is
\(\pi=\alpha^{-1}(\bar\pi)\). Its block-state tensor,
density, shift and residue faces agree with the
\(\bar\pi\)-term. Closed direct image and open
base change give the exceptional equivalence.

Saturation for two collisions is saturation for
their composite, so the same identifications give
all the required coherences.

On the cross-disjoint open of a consecutive cut
\(I=A|D\), a surviving partition is uniquely
\(\pi_A\sqcup\pi_D\). Its states and densities
split into the two consecutive blocks. Localization
and this equivalence define the cut coproduct.

For any further cuts all composites use the same
disjoint open, partition decomposition and ordered
tensor factors. Thus the coproduct is coassociative
and commutes with collisions. Empty cuts use the
specified coefficient \(k\). There are at most
\(|I|\) positive output blocks, which proves
conilpotence in each arity.
\end{proof}

At three positions put
\(H_X=\omega_X\otimes_k H_\kappa\), and let each
\(L\) denote the localized two-block coefficient
with two \(H_\kappa\)-factors. Writing \(j\) for
the three-point configuration open, the terms are
\[
 \begin{aligned}
 (C_\kappa)^{-3}&=j_*j^*(H_X^{\boxtimes3}),\\
 (C_\kappa)^{-2}&=
    i_{12,*}L_{12|3}\oplus i_{23,*}L_{1|23}
                              \oplus i_{13,*}L_{13|2},\\
 (C_\kappa)^{-1}&=i_{123,*}H_X.
 \end{aligned}
\]

The third summand is present even though the cut tensor uses consecutive
cuts. For independent currents \(J,P\), the outer
contraction in \(J(z_1)P(z_2)J(z_3)\) is
\[
                  \frac{\kappa P(z_2)}{(z_1-z_3)^2}.
\]
It survives on \(13|2\). Both adjacent contractions
vanish. The scalar vacuum expectation of this term
is zero. Thus that expectation does not retain the
supported operation.

Insert \(\mathbf1\) in every block of the scalar
residue coalgebra \(C^{\mathrm{run}}(k)\) from the
preceding chapter. Since \(Y(\mathbf1,u)\mathbf1=\mathbf1\),
these maps commute with every face, collision and cut.
They give a coalgebra map
\begin{equation}\label{v4-eq:singular-vacuum-map}
                    C^{\mathrm{run}}(k)\longrightarrow C_\kappa.
\end{equation}
Its singleton is \(\omega_X[1]\to
\omega_X\otimes H_\kappa[1]\), induced by the local
vacuum. The source is not the isolated empty unit \(K\).

\begin{definition}\label{v4-def:singular-background}
The \emph{Heisenberg geometric background} on the word
carrier is the augmented associative algebra
\[
                 A_\kappa=\Omega_{\mathrm{cut}}C_\kappa.
\]
The vacuum background is
\(A_{\mathrm{vac}}=\Omega_{\mathrm{cut}}C^{\mathrm{run}}(k)\).
\end{definition}

\begin{corollary}\label{v4-cor:singular-background-comparison}
There is an associative map
\(A_{\mathrm{vac}}\to A_\kappa\), and
\[
 B_{\mathrm{cut}}A_\kappa\simeq C_\kappa,
 \qquad (A_\kappa)_\varnothing=k.
\]
The framed two-point boundary of this bar recovers
every mode of \(H_\kappa\).
\end{corollary}

\begin{proof}
Apply cut cobar to \eqref{v4-eq:singular-vacuum-map}.
The connected coalgebra satisfies the positive-cut
hypotheses of Theorem~\ref{v4-thm:bar-cut-duality}.
That theorem gives the bar comparison. Every positive
cobar word vanishes at the empty input, so only its
unit remains there.

In the framed two-point boundary,
\[
 r(f;a,b)=\operatorname{pp}_u(fY(a,u)b),
 \qquad [e_0]r(u^n;a,b)=a_{(n)}b.
\]
All integer Laurent tests are allowed. This proves
the mode recovery, including the central first product.
\end{proof}

The augmentation \(A_\kappa\to K\) is a geometric
augmentation of a cut cobar. It is not a unital
vertex-algebra map \(H_\kappa\to k\), which
Proposition~\ref{v4-prop:singular-no-augmentation} excludes.
Likewise cobar multiplication is concatenation. The
field products are recovered from the framed contacts.
The three-channel boundary identity does not identify
those contacts with strictly associative composition
of two consecutive residues.

For a unital differential coefficient algebra \(B\),
one may instead retain both the local field factor
and the runs of \(B\). The partition coefficient is
\[
 \omega_{X^\pi}\otimes H_\kappa^{\otimes|\pi|}
                         \otimes B^{\otimes\ell(\pi)}.
\]
Its merger combines \(Y_{H_\kappa}\) with the ordered
run map and Taylor factor of
\eqref{v4-eq:regular-run-face}. The same three-channel
proof applies: every route has the same ordered
\(B\)-coefficient after its formal translations.
This gives \(C^{\mathrm{run}}(H_\kappa,B)\) and its
cut cobar, with empty coefficient \(k\).

The relative envelope of
Chapter~\ref{v4-ch:relative-coefficient-envelopes} is
\(A_\kappa\otimes_k B\), with empty coefficient \(B\).
These are distinct constructions. An equivalence
between them would already imply \(B\simeq k\)
after evaluation at the empty word. For an ordinary
algebra of dimension greater than one this is
impossible. A proposed comparison must therefore
specify its relative base and its effect on the
field contacts. The notation of scalar extension
alone supplies neither.

\section{The derived coefficient of the local vacuum}
\label{v4-sec:singular-vacuum-width}

There is a second geometric vacuum construction.
The dualizing coalgebra \(J\) of
Proposition~\ref{v4-prop:word-vacuum-coalgebra}
has \(J_n=\omega_{X^n}[n]\) in positive arity.
Its cut coproducts are localization maps.
Write
\[
             Q_n=(\Omega_{\mathrm{cut}}J)_n,\qquad n>0.
\]
This definition uses \(J\) itself.
Identifying it with the scalar residue background
would require a comparison with
\(C^{\mathrm{run}}(k)\) preserving the coalgebra
and collision maps.

We compute \(Q_n\), including its support extensions.
Let \(X\) now be any smooth separated
geometrically connected curve over \(k\), and
put \(Y_n=X^n\).
Use the closed support \(Z_n\) of
\eqref{v4-eq:word-collision-support}.
Let \(U_t\) be the open where every
position to the left of cut \(t\)
differs from every position to its right.
For \(S\subseteq\{1,\ldots,n-1\}\), put
\(U_S=\bigcap_{t\in S}U_t\), with inclusion
\(j_S\), and \(U_\varnothing=Y_n\).

Define the localization cube
\begin{equation}\label{v4-eq:singular-vacuum-cube}
 \mathsf C_n^r=\bigoplus_{|S|=r}j_{S,*}\omega_{U_S},
 \qquad
 d_{S,S\cup\{t\}}
       =(-1)^{|\{s\in S:s<t\}|}\mathrm{loc}.
\end{equation}
All derivatives normal to a diagonal are retained.
Each localization permits arbitrary finite pole
order. No completion in arity or pole order
is used.

\begin{proposition}\label{v4-prop:singular-vacuum-coefficient}
For \(n\geq2\), there are natural comparisons
\[
 \mathsf C_n\simeq R\Gamma_{Z_n}\omega_{Y_n},
 \qquad Q_n\simeq\mathsf C_n[n-1].
\]
They retain the cube and its augmentation.
With ordered density wedges and cobar
desuspensions, the second comparison is
the vertexwise sign \((-1)^{n|S|}\).
\end{proposition}

\begin{proof}
The opens \(U_t\) cover \(Y_n\setminus Z_n\).
The cube including its empty vertex computes
the fiber of restriction to this union.
Every intersection is locally a principal
divisor complement, so its localization is exact.
This proves the first comparison.

A set \(S\) gives \(|S|+1\) consecutive
cobar blocks. Their density shifts sum
to \([n]\), while desuspending every block
contributes \([-|S|-1]\).
Thus its coefficient lies in degree
\(|S|+1-n\), as in \(\mathsf C_n[n-1]\).
Fix the suspension identification
\[
 (tc_1)\otimes\cdots\otimes(tc_b)
   \longmapsto
 (-1)^{\sum_{a=1}^b(b-a)|c_a|}
                 (c_1\otimes\cdots\otimes c_b)[-b],
 \qquad t=[-1].
\]

Splitting a block at the new cut \(t\)
has sign \((-1)^{1+a}\), where
\(a=|\{s\in S:s<t\}|\).
The shifted cube has sign \((-1)^{n-1+a}\).
The identity
\[
 (-1)^{n(|S|+1)}(-1)^{1+a}
             =(-1)^{n-1+a}(-1)^{n|S|}
\]
proves the claimed chain isomorphism.
Both maps use the same localization at
every face, including the empty vertex.
\end{proof}

\subsection{An order with divisor intersections}

For partitions \(p,q\) of \([n]\), write
\(p\leq q\) when every \(p\)-block
lies in a \(q\)-block. Their join is
\(p\vee q\), and \(c(p)=n-|p|\).
Then
\(\Delta_p\cap\Delta_q=\Delta_{p\vee q}\).
Call a partition indecomposable when
its blocks cross every consecutive cut.

By Theorem~\ref{v4-thm:word-support-components},
the minimal indecomposable partitions are
the crossing chains
\[
 1=i_1<i_2<j_1<i_3<j_2<\cdots
                         <i_c<j_{c-1}<j_c=n,
\]
whose nonsingleton blocks are \((i_a,j_a)\).
The one-pair chain is \((1,n)\).
Let \(\mathcal A_n\) be their set,
ordered lexicographically by
\begin{equation}\label{v4-eq:singular-vacuum-order}
              (-j_1,i_1,-j_2,i_2,\ldots,-j_c,i_c).
\end{equation}
Write \(N\prec M\) for this order.
No such list properly extends another:
the last right endpoint is already \(n\).

For an indecomposable \(p\), replace each
nonsingleton block by its hull
\((\min B,\max B)\).
Start with the hull containing \(1\).
At a current right endpoint \(r<n\),
choose the eligible hull with largest right
endpoint, where eligibility means
\[
                         \min B<r<\max B.
\]
Continue to \(n\), and leave unused
positions singleton. Denote this refinement
by \(g(p)\).

\begin{lemma}\label{v4-lem:singular-greedy-exchange}
The partition \(g(p)\) is the least
crossing chain refining \(p\).
For \(N\prec M\), let \(e\) be the
first pair of \(N\) that differs from \(M\),
and put \(p=M\vee e\).
Then
\begin{equation}\label{v4-eq:singular-greedy-exchange}
 g(p)\prec M,\qquad M\vee g(p)=p\leq M\vee N,
 \qquad c(p)=c(M)+1.
\end{equation}
\end{lemma}

\begin{proof}
The cut at \(r\) is crossed by a
block with \(\min B\leq r<\max B\).
Equality on the left is impossible:
\(r\) already belongs to the previously
chosen block, whose maximum is \(r\).
Thus an eligible hull exists.
Distinct blocks have distinct maxima.
The right endpoint strictly increases,
so the process terminates at \(n\).

Write its selected pairs as \((i_a,j_a)\).
For \(a\geq3\), one has \(i_a>j_{a-2}\).
Equality would share an endpoint.
A smaller value would make this hull
eligible one step earlier, with a larger
right endpoint than the greedy choice.
Together with \(i_a<j_{a-1}<j_a\),
these inequalities give the crossing-chain order.

Compare the output with any chain refining \(p\).
At the first differing pair, the greedy
right endpoint is at least the competing one.
For the first pair this follows from
the hull of the block containing \(1\).
For a later pair, the competing pair
lies in an eligible hull.
If the right endpoints agree, they
belong to the same block, and the
greedy left endpoint is its minimum.
This proves lexicographic minimality.

Now let the first differing pair be
\(e=(i,j)\), and let the corresponding
pair of \(M\) be \((i_a,j_a)\).
Since \(N\prec M\), either \(j>j_a\),
or \(j=j_a\) and \(i<i_a\).
Suppose \(g(M\vee e)=M\).
In the first position, the new block
containing \(1\) has maximum larger than \(j_1\),
contradicting this supposition.

At a later position put \(r=j_{a-1}\).
The block containing \(e\) has
\(\min B\leq i<r<j\leq\max B\).
It is eligible.

A larger right endpoint contradicts greedy
maximality. If its maximum equals \(j_a\),
uniqueness of that maximum forces this block
to be chosen, with left endpoint at most
\(i<i_a\). This is again a contradiction.
Hence \(g(p)\prec M\).

Two distinct minimal indecomposable partitions
are incomparable. Therefore
\[
                     M<M\vee g(p)\leq p.
\]
Adding one pair raises partition rank
by at most one.
The strict inequality forces
\(M\vee g(p)=p\) and \(c(p)=c(M)+1\).
Since \(e\) is a pair of \(N\),
one also has \(p\leq M\vee N\).
\end{proof}

For \(M\in\mathcal A_n\), define
\[
 \mathscr D_M=\{\sigma:\sigma\text{ covers }M
                              \text{ and }g(\sigma)\prec M\}.
\]
A cover merges exactly two \(M\)-blocks.
The exchange lemma proves the attachment equation
\begin{equation}\label{v4-eq:singular-vacuum-attachment}
 D_M:=\Delta_M\cap\bigcup_{N\prec M}\Delta_N
                         =\bigcup_{\sigma\in\mathscr D_M}\Delta_\sigma.
\end{equation}
Indeed every earlier intersection lies
in a cover supplied by the lemma.
Conversely \(\Delta_\sigma\) lies in both
\(\Delta_M\) and the earlier \(\Delta_{g(\sigma)}\).

Each equality of two distinct block
coordinates is a Cartier divisor on
the smooth product \(\Delta_M\).
Consequently \(D_M\) is empty or is
the support of an effective Cartier divisor.
On the affine line its equation is
\[
 h_M=\prod_{\substack{B<B'\\g(M/(B\sim B'))\prec M}}
                  (y_B-y_{B'}),\qquad D_M=V(h_M).
\]
Blocks are ordered by their least label.
The empty product is \(1\).

\subsection{The support maps and exact width}

Suppose \(i:M\hookrightarrow Y\) is a closed embedding
of smooth varieties of codimension \(c\), and \(A\subset Y\)
is closed with \(M\not\subset A\).
If \(D=A\cap M\) is a divisor support
or is empty, there is a natural triangle
\begin{equation}\label{v4-eq:singular-relative-support}
 R\Gamma_A\omega_Y\longrightarrow R\Gamma_{A\cup M}\omega_Y
       \longrightarrow i_+\omega_M(*D)[-c]
       \longrightarrow R\Gamma_A\omega_Y[1].
\end{equation}
Here \(i_+\) is closed direct image
with all finite normal derivatives.

To construct it, use the support square
\[
 \begin{matrix}
 R\Gamma_D\omega_Y&\longrightarrow&R\Gamma_M\omega_Y\\
 \downarrow&&\downarrow\\
 R\Gamma_A\omega_Y&\longrightarrow&R\Gamma_{A\cup M}\omega_Y .
 \end{matrix}
\]
Mayer--Vietoris makes this a homotopy pushout.
Smooth purity identifies the upper row with
\[
 i_+R\Gamma_D^M\omega_M[-c]\longrightarrow i_+\omega_M[-c].
\]
Its cofiber is \(i_+\omega_M(*D)[-c]\).
Equality of the horizontal cofibers gives
\eqref{v4-eq:singular-relative-support}, including
its connecting map.

The normalization can be checked directly.
For ordered normal parameters \(u_1,\ldots,u_c\),
the two-term localization complexes have
cohomology only in degree \(c\), with
\[
 \left[\frac{du_1\wedge\cdots\wedge du_c\wedge\omega_M}
                    {u_1^{a_1+1}\cdots u_c^{a_c+1}}\right]
 \longmapsto
       \frac{\omega_M\otimes\partial_{u_1}^{a_1}\cdots
                             \partial_{u_c}^{a_c}}
                 {a_1!\cdots a_c!}.
\]
Right differentiation of densities gives
the positive factorials.
Normal-coordinate changes transform the
density and residue determinant together,
so these formulas glue.

The width of a nonzero bounded complex
is the difference between its extreme
nonzero cohomological degrees.

\begin{theorem}\label{v4-thm:singular-vacuum-width}
For every \(n\geq2\),
\[
 \begin{aligned}
 H^j(Q_n)&=0
       &&\text{unless }2-n\leq j\leq\lfloor n/2\rfloor+1-n,\\
 H^{c+1-n}(Q_n)&\ne0
       &&(1\leq c\leq\lfloor n/2\rfloor).
 \end{aligned}
\]
Thus \(\operatorname{width}(Q_n)=\lfloor n/2\rfloor-1\).
The successive support maps in the order
\eqref{v4-eq:singular-vacuum-order} have cofibers
\[
                     i_{M,+}\omega_{\Delta_M}(*D_M)[-c(M)]
\]
before the common shift \([n-1]\).
\end{theorem}

\begin{proof}
Attach the components in their specified order.
Each has codimension at most
\(m=\lfloor n/2\rfloor\).
Equation~\eqref{v4-eq:singular-vacuum-attachment}
and triangle~\eqref{v4-eq:singular-relative-support}
construct the asserted tower.
Its factors have cohomology only in
degrees between \(1\) and \(m\).
The long exact sequences propagate this bound.
Proposition~\ref{v4-prop:singular-vacuum-coefficient}
then gives the vanishing range after the shift.

For each \(1\leq c\leq m\), the component
count in Theorem~\ref{v4-thm:word-support-components}
supplies a chain \(M\) of length \(c\).
Choose a geometric point with equality
partition exactly \(M\).
No other minimal indecomposable partition
occurs there. Locally \(Z_n\) is therefore
this smooth codimension-\(c\) diagonal.
Its class
\[
                      \left[\frac{\omega_{Y_n}}{u_1\cdots u_c}\right]
\]
is nonzero in the normal Čech quotient.

Faithful extension to a geometric field
detects nonzero sheaves over \(k\).
Thus every asserted degree occurs.
\end{proof}

The upper sheaf bound also holds for
any quasicoherent coefficient over any field.
At a codimension-\(c\) component, its support
has \(c\) local equations.
Its intersection with the preceding union
has those equations and one lifted divisor
equation. Their Čech complexes vanish
above \(c\) and \(c+1\), respectively.
The support Mayer--Vietoris sequence therefore
propagates vanishing above \(\lfloor n/2\rfloor\).
This argument uses no characteristic-zero purity.

Shifts preserve width.
Consequently no arity-dependent shifts can put
all \(Q_n\) into a fixed finite range
of cohomological degrees.
For \(n\geq4\), even a shifted heart
object cannot represent \(Q_n\).

\subsection{The four-position extension}

Take \(X=\mathbb A^1_k\) and coordinates
\[
       a=z_1-z_4,\quad b=z_1-z_3,\quad
       c=z_2-z_4,\quad s=z_4.
\]
Then \(Z_4=V(a)\cup V(b,c)\).
The ordered attachments first add \(14\),
then \(13|24\) along its divisor \(a=0\).
Put \(E_x=k[x,x^{-1}]/k[x]\).

Suppressing the common density and \(k[s]\)-factor,
triangle~\eqref{v4-eq:singular-relative-support} gives
\[
 H^1_{Z_4}=E_a\otimes k[b,c],\qquad
 H^2_{Z_4}=k[a,a^{-1}]\otimes E_b\otimes E_c.
\]
All other local cohomology vanishes.
In particular there is an exact sequence
\begin{equation}\label{v4-eq:singular-four-extension}
 0\longrightarrow k[a]\otimes E_b\otimes E_c
 \longrightarrow H^2_{Z_4}
 \longrightarrow E_a\otimes E_b\otimes E_c
 \longrightarrow0.
\end{equation}

This sequence does not split even as
a sequence of ordinary modules.
The quotient has a nonzero element killed
by \(a\), whereas multiplication by \(a\)
is invertible on the middle module.
A module section cannot preserve that relation.
After the shift \([3]\), the cohomology
of \(Q_4\) lies in degrees \(-2,-1\).

The tower retains its support and connecting
maps. Its ordered factors cannot replace
the coefficient by a direct sum.
The ordering of components is auxiliary.
Comparisons of these particular towers under
cuts and monotone collisions must additionally
preserve their extension maps and recover
the operations of the unfiltered coefficient.

\section{Infinitely many species at fixed energy}
\label{v4-sec:singular-hilbert-fields}

For the analytic constructions set \(k=\mathbb C\)
and normalize the level to one. Let \(H\) be a real
separable Hilbert space. Its complexification
\(H_{\mathbb C}\) has both the Hermitian norm and
the complex bilinear extension \((\ ,\ )\) of the
real inner product. The bilinear pairing supplies
contractions. The Hermitian norm supplies estimates.
For \(m\geq1\), let \(H_m\) be its copy of energy
\(m\), with squared norm \(m\|v\|^2\).

Define the Hilbert energy sector and the algebraic
energy sum by
\begin{equation}\label{v4-eq:singular-hilbert-sectors}
 \mathcal V_E=
    \bigoplus_{\sum m r_m=E}
       \widehat{\bigotimes}_{m\geq1}
              \widehat{\operatorname{Sym}}^{r_m}H_m,
 \qquad
 \mathcal V=\bigoplus_{E\geq0}\mathcal V_E.
\end{equation}
The symmetric tensors have the Fock normalization:
a product of \(r\) identical unit creators has
squared norm \(r!\). Each displayed sum at fixed
\(E\) has finitely many occupancy patterns. Only
the sum over all energies remains algebraic.
Denote its Hilbert completion by
\(\widehat{\mathcal V}\).

The creation and annihilation modes satisfy
\[
 [a_m(v),a_n(w)]=m\delta_{m+n,0}(v,w),
 \qquad a_0(v)=0,
 \qquad a_m(v)\mathbf1=0\quad(m>0).
\]
They are complex linear in the species vector.
For real vectors the paired creation and annihilation
operators are adjoint. On finite species tensors set
\[
 J_v(z)=\sum_m a_m(v)z^{-m-1},
 \qquad
 Y\left(\prod_{j=1}^p a_{-m_j}(v_j)\mathbf1,z\right)
   =:\prod_{j=1}^p
         \frac{\partial_z^{m_j-1}J_{v_j}(z)}{(m_j-1)!}:.
\]
We next extend each coefficient of this formula.

\begin{lemma}\label{v4-lem:singular-cross-contractions}
Contract \(r\) specified legs of a Hilbert tensor
\(a\in H_{\mathbb C}^{\widehat\otimes p}\) against
\(r\) specified legs of
\(b\in H_{\mathbb C}^{\widehat\otimes q}\).
The resulting tensor has norm at most
\(\|a\|\|b\|\). Consequently each mode extends
to a bounded bilinear map
\begin{equation}\label{v4-eq:singular-mode-bound}
 \mathcal V_E\times\mathcal V_F
   \longrightarrow\mathcal V_{E+F-n-1},
 \qquad (a,b)\longmapsto a_{(n)}b.
\end{equation}
A negative-energy sector is understood to be zero.
\end{lemma}

\begin{proof}
In a real orthonormal basis group the contracted
indices together. A coefficient of the contraction
has the form \(\sum_\gamma a_{\alpha\gamma}
b_{\gamma\beta}\). Cauchy--Schwarz gives
\[
 \sum_{\alpha,\beta}
     \left|\sum_\gamma a_{\alpha\gamma}b_{\gamma\beta}\right|^2
 \leq
 \left(\sum_{\alpha,\gamma}|a_{\alpha\gamma}|^2\right)
 \left(\sum_{\gamma,\beta}|b_{\gamma\beta}|^2\right).
\]
The same inequality applies to the complex bilinear
contraction. Symmetrization is an orthogonal projection.

At fixed \(E,F,n\), the Wick formula has finitely
many occupancy, pairing and Taylor patterns. Their
output energy is \(E+F-n-1\), so no unbounded
choice of remaining creation modes occurs. The
mode weights and Fock factorials contribute only
a finite constant for these fixed indices.
Summing the bounds proves the stated extension.
\end{proof}

To linearize a bounded bilinear coefficient,
we use the completed \emph{projective} tensor
product of its two input sectors.
No bounded linear map on their Hilbert
tensor product is inferred from bilinear boundedness.

\begin{theorem}\label{v4-thm:singular-hilbert-vertex}
The extended coefficient maps make \(\mathcal V\)
a local vertex algebra. Its vacuum is \(\mathbf1\),
and its translation satisfies
\[
 T\mathbf1=0,\qquad
 [T,a_{-m}(v)]=m a_{-m-1}(v).
\]
For energies \(E,F\), the locality exponent \(E+F\)
is independent of a species cutoff.
\end{theorem}

\begin{proof}
Output energy proves lower truncation in
\eqref{v4-eq:singular-mode-bound}. Translation changes
only finitely many occupancy patterns on
\(\mathcal V_E\), and is bounded into
\(\mathcal V_{E+1}\). Its identities and the vacuum
identities therefore extend from finite species.

In a Wick contraction the pole order of a pair
of derivative currents is the sum of their
energies. Thus \(E+F\) bounds the total pole
order before taking any species limit.
Clear these poles and apply locality to a
third fixed-energy state. Each coefficient
is a finite sum of bounded compositions
of mode maps. It vanishes on finite species
and hence on their dense completion. The mode
Jacobi identities extend in the same way: the
energy bound makes each of their fixed-index
sums finite, uniformly in the approximants.
\end{proof}

There is a precise restriction in this argument.
Contraction between two tensors is bounded.
Contraction of two legs of one arbitrary Hilbert
tensor need not be. For an orthogonal projection
\(P_N\) of rank \(N\), its tensor norm is
\(\sqrt N\), while its trace is \(N\).
Vacuum normal ordering has removed such internal
contractions from the fixed-coordinate definition.
They will return under general coordinate changes.

\subsection{Quadratic fields and all their poles}

For a real symmetric Hilbert--Schmidt operator
\(A\) on \(H\), define
\begin{equation}\label{v4-eq:singular-quadratic-state}
 q_A=\sum_{i,j}(e_i,Ae_j)
        a_{-1}(e_i)a_{-1}(e_j)\mathbf1,
 \qquad W_A(z)=Y(q_A,z).
\end{equation}
The Hilbert--Schmidt norm is
\(\|A\|_2^2=\sum_{i,j}|(e_i,Ae_j)|^2\).

\begin{proposition}\label{v4-prop:singular-quadratic-isometry}
The state \(q_A\in\mathcal V_2\) is independent
of the real orthonormal basis and satisfies
\[
                         \|q_A\|^2=2\|A\|_2^2.
\]
Conversely, for bounded symmetric \(A\), convergence
of \(q_{P_NAP_N}\) for finite-rank orthogonal
\(P_N\to1\) strongly implies that \(A\) is
Hilbert--Schmidt.
\end{proposition}

\begin{proof}
The two pairings of two creation factors give
the norm formula in finite rank. The intrinsic
identification of symmetric Hilbert--Schmidt
operators with symmetric Hilbert tensors then
extends it and proves basis independence.
For the converse the norm formula makes the
compressions Cauchy in Hilbert--Schmidt norm.
Their limit agrees with \(A\) in every matrix
element, by strong convergence of \(P_N\).
\end{proof}

For symmetric \(A,B\), set
\[
 S_{AB}=\tfrac12(AB+BA),\qquad
 D_{AB}(w)=\sum_{i,j}(e_i,ABe_j)
                      :\partial J_{e_i}(w)J_{e_j}(w):.
\]
The product \(AB\) is trace class, hence
Hilbert--Schmidt, so the latter is the field
of an energy-three state.

\begin{theorem}\label{v4-thm:singular-quadratic-OPE}
The full singular quadratic product is
\begin{equation}\label{v4-eq:singular-quadratic-OPE}
 W_A(z)W_B(w)\sim
    \frac{2\operatorname{Tr}(AB)}{(z-w)^4}
   +\frac{4W_{S_{AB}}(w)}{(z-w)^2}
   +\frac{4D_{AB}(w)}{z-w}.
\end{equation}
Its coefficients converge in their Hilbert
energy sectors under Hilbert--Schmidt approximation.
\end{theorem}

\begin{proof}
The two double contractions give the fourth-order
pole. The four single contractions give
\[
 \frac4{(z-w)^2}\sum_{i,j}(e_i,ABe_j)
                          :J_{e_i}(z)J_{e_j}(w):.
\]
Its constant Taylor term is \(4W_{S_{AB}}\),
since the equal-position creation factors are
symmetric. Its linear term is \(4D_{AB}\).
All remaining Taylor terms are regular.
The inequality \(\|AB\|_1\leq\|A\|_2\|B\|_2\)
controls both the trace and the coefficient
tensors, and passes the finite-species calculation
to the limit.
\end{proof}

The simple pole retains the order of the two
operators. Indeed
\[
 D_{AB}+D_{BA}=\partial W_{S_{AB}},\qquad
 4D_{AB}=2\partial W_{S_{AB}}+2(D_{AB}-D_{BA}).
\]
For \(A=\operatorname{diag}(1,0)\) and
\(B=\left(\begin{smallmatrix}0&1\\1&0\end{smallmatrix}\right)\),
the last term is
\[
                  2(:\partial J_1J_2:-:\partial J_2J_1:).
\]
Its two state components have modes \((-2,-1)\),
are orthogonal and each have squared norm two.
The total squared norm is sixteen. The Jordan
product alone would miss this nonzero field.

For orthogonal rank-one projections \(P_1,P_2\),
the outer contraction in
\(W_{P_1}(z_1)W_{P_2}(z_2)W_{P_1}(z_3)\) is
\begin{equation}\label{v4-eq:singular-quadratic-outer}
                   \frac{2W_{P_2}(z_2)}{(z_1-z_3)^4}.
\end{equation}
All contractions meeting the middle species vanish.
The normal test \((z_1-z_3)^3\) therefore gives
the nonzero residue \(2W_{P_2}\), with squared
state norm eight. Its scalar vacuum expectation
is zero. This is the quadratic version of the
supported operation lost by scalar evaluation.

On the specified affine coordinate, the right
module \(\mathcal M=\omega_X\otimes\mathcal V\)
has derivative \(-\partial_z-T\). Its local
chiral operation is
\begin{equation}\label{v4-eq:singular-hilbert-chiral}
 j_*j^*(\mathcal M\boxtimes\mathcal M)
      \longrightarrow\Delta_!\mathcal M,
 \qquad f(u,z)(a\boxtimes b)
      \longmapsto\operatorname{pp}_u(f(u,z)Y(a,u)b).
\end{equation}
Here the external tensor is algebraic when
\(\mathcal V\) is regarded as a vector space.
Fixed-energy continuous coefficients linearize
on projective external tensors.

Lower truncation makes every displayed normal
sum finite. Translation proves differential
linearity. Skew symmetry and Jacobi prove the
chiral identities in the full normal module.
This construction uses its stated auxiliary
curve. It does not identify that curve with
the support of an arithmetic random field.

\subsection{The center and the energy grading}

The ordinary vertex center consists of states
whose field has no singular product with any
state. Skew symmetry gives the equivalent
condition that every nonnegative mode of every
state annihilates the central state.

\begin{proposition}\label{v4-prop:singular-ordinary-centre}
The ordinary vertex center of \(\mathcal V\)
is \(\mathbb C\mathbf1\).
\end{proposition}

\begin{proof}
A central state is killed by every positive
current mode \(a_m(h)\). At fixed occupancy,
these modes contract each occupied leg against
\(h\). The nondegenerate Hilbert pairing makes
a tensor with all such contractions zero equal
to zero. Different occupancy patterns remain
distinguishable after the contracted mode is
fixed. All positive-energy components therefore
vanish. The vacuum is central.
\end{proof}

\begin{proposition}\label{v4-prop:singular-conformal-boundary}
For a finite-rank orthogonal projection \(P\),
\(L_P=\tfrac12W_P\) satisfies the Virasoro
operator product with central charge
\(\operatorname{rank}P\). If \(H\) is infinite
dimensional, no state of \(\mathcal V\) is a
conformal vector implementing its energy grading.
\end{proposition}

\begin{proof}
Since \(P^2=P\),
\eqref{v4-eq:singular-quadratic-OPE} becomes
\[
 L_P(z)L_P(w)\sim
  \frac{\operatorname{rank}(P)/2}{(z-w)^4}
       +\frac{2L_P(w)}{(z-w)^2}
       +\frac{\partial L_P(w)}{z-w}.
\]
Its grading concerns the projected species.
The state has squared norm
\(\operatorname{rank}(P)/2\).

A conformal vector \(\omega\) implementing the
given energy must have energy two, because
\(L_0\omega=2\omega\). Every vector there is
\(q_A+a_{-2}(v)\mathbf1\), for a complex
symmetric Hilbert--Schmidt tensor \(A\).
Its mode \((1)\) acts on the energy-one species
space by \(2A\). The improvement field
\(\partial J_v\) contributes
\(-a_0(v)=0\) to that mode.
The required identity action would force
\(2A=1_H\), which is not Hilbert--Schmidt
in infinite dimension.
\end{proof}

Thus vacuum and translation survive without an
energy-grading conformal state. The scalar
ordinary center just computed also leaves the
derived center and its geometric actions as
separate objects of comparison.

\subsection{A prime-labeled current cover}

For each prime \(p\), take two orthogonal
level-one currents \(a^{(p)},b^{(p)}\).
Their rank-two vertex algebra
\(\mathcal H^{(p)}\) has conformal state
\[
 \omega_p=\frac12\left(
        (a^{(p)}_{-1})^2+(b^{(p)}_{-1})^2\right)\mathbf1_p.
\]
The preceding calculation gives central
charge two and integer oscillator
grading \(L_0^{(p)}\).
For finite sets of primes \(S\), the
maps
\[
 \bigotimes_{p\in S}\mathcal H^{(p)}
       \longrightarrow
 \bigotimes_{p\in S'}\mathcal H^{(p)},\qquad S\subset S',
\]
insert the vacuum in each new factor.
They preserve all field products and
translation. Their filtered union defines
the algebraic vacuum-restricted product
\[
                \mathcal H_{\mathbb P}
                     =\bigotimes_p'\mathcal H^{(p)}.
\]

Every state uses finitely many primes.
For a finite set \(S\), its conformal
state is \(\sum_{p\in S}\omega_p\),
of central charge \(2|S|\).
The infinite sum is not a state of
\(\mathcal H_{\mathbb P}\).

There is also the logarithmically weighted
operator
\[
                  H_{\mathrm{ar}}
                     =\sum_p(\log p)L_0^{(p)}
\]
on this algebraic product. Each state
makes this sum finite.
It assigns energy \(m\log p\) to
a mode-\(m\) creator at \(p\).
Its weights differ from the oscillator
grading. Rescaling \(\omega_p\) does
not preserve its Virasoro normalization.

With species space
\(H=\ell^2(\mathbb P;\mathbb R^2)\),
the algebraic product is dense in each
Hilbert energy sector constructed above.
An extension of \(H_{\mathrm{ar}}\)
to a closed Hilbert operator requires
its own domain.
A supported map from this current cover
to independently defined arithmetic
coefficients is further data.

\section{Coordinate changes and coefficient norms}
\label{v4-sec:singular-coordinates}

A coordinate change replaces the kernel subtracted
in a normally ordered field. For a quadratic field
this replacement contracts its two coefficient
indices. The resulting trace determines whether
the change acts on the chosen completion.
Assume throughout this and the following sections
that \(H\) is infinite dimensional.

Let \(f\) be a locally invertible holomorphic germ.
The difference of the two current kernels is
\begin{equation}\label{v4-eq:singular-coordinate-kernel}
 H_f(z,u)=\frac{f'(z)f'(u)}{(f(z)-f(u))^2}
                         -\frac1{(z-u)^2}.
\end{equation}
It is regular along the diagonal. The Taylor
expansion of \(f(u+s)\) through order three gives
\begin{equation}\label{v4-eq:singular-Schwarzian}
 H_f(z,z)=\tfrac16\{f,z\},\qquad
 \{f,z\}=\frac{f'''(z)}{f'(z)}
              -\frac32\left(\frac{f''(z)}{f'(z)}\right)^2.
\end{equation}

Inserting the intermediate kernel under a second
coordinate change gives the exact cocycle
\begin{equation}\label{v4-eq:singular-kernel-cocycle}
 H_{f\circ g}(z,u)
       =H_g(z,u)+g'(z)g'(u)H_f(g(z),g(u)).
\end{equation}
Consequently a finite-rank quadratic coefficient
has the transformation law
\begin{equation}\label{v4-eq:singular-coordinate-quadratic}
 W_A^{(z)}(z)
   =f'(z)^2W_A^{(w)}(f(z))
          +\tfrac16\operatorname{Tr}(A)\{f,z\}\mathbf1,
 \qquad w=f(z).
\end{equation}
The sign follows by subtracting the old vacuum
kernel from the pullback of the new one.

For a germ fixing zero, put
\[
 B_f(m,n)=[z^{m-1}u^{n-1}]H_f(z,u),\qquad m,n\geq1.
\]
The coordinate map on a finite-species state sums
over its internal pairings. A paired pair of
legs in modes \(m,n\) contributes its species
pairing times \(B_f(m,n)\). The unpaired legs
undergo the finite triangular mode substitution
defined below.

\begin{theorem}\label{v4-thm:singular-coordinate-obstruction}
If some \(B_f(m,n)\ne0\), this coordinate map
on finite species is not closable in the
Hilbert energy spaces of
\eqref{v4-eq:singular-hilbert-sectors}.
Precisely the fractional-linear germs fixing
zero act continuously on every bounded-energy
space with the specified current and vacuum.
\end{theorem}

\begin{proof}
For an orthonormal sequence \((e_i)\), consider
\[
 u_N^{mn}=\frac1N\sum_{i=1}^N
                a_{-m}(e_i)a_{-n}(e_i)\mathbf1.
\]
Its squared norm is \(2m^2/N\) if \(m=n\)
and \(mn/N\) otherwise. Thus \(u_N^{mn}\to0\).
Its uncontracted coordinate image also tends
to zero, because only finitely many bounded
mode substitutions occur. Its contracted
image is the fixed scalar \(B_f(m,n)\mathbf1\).
The closure of the graph therefore contains
a nonzero vector over zero.

All coefficients \(B_f(m,n)\) vanish exactly
when \(H_f\) vanishes as a germ. In that case
\(\{f,z\}=0\). Put \(v=f''/f'\). Its equation
is \(v'=v^2/2\). Integration gives a
fractional-linear function, including the
affine case.

Conversely direct substitution
of \(f(z)=(az+b)/(cz+d)\), \(ad-bc\ne0\),
gives \(H_f=0\). Its coordinate map and inverse
then consist only of bounded mode substitutions
on every finite set of energy sectors.
\end{proof}

The germ \(f(z)=z+z^4\) has
\(B_f(1,1)=0\) but \(B_f(1,2)=2\).
Testing only the Schwarzian at the marked
point therefore misses an obstruction at
higher energy. Restricting to traceless
finite-rank quadratics also fails to give a
field algebra: for \(A=\operatorname{diag}(1,-1)\),
the quadratic product contains \(A^2\), whose
trace is two.

\subsection{Projective tensors and all coordinate jets}

For a finite tensor define its projective norm by
\[
 \|u\|_\pi=
 \inf_{u=\sum_\alpha v_{\alpha,1}\otimes\cdots\otimes v_{\alpha,p}}
       \sum_\alpha\prod_j\|v_{\alpha,j}\|.
\]
Write \(\widehat\otimes_\pi\) for the completion
and \(\widehat{\operatorname{Sym}}_\pi^r\) for
the image of its continuous symmetrizing
projection. Use the same Fock factorial weights
as before and set
\begin{equation}\label{v4-eq:singular-projective-sectors}
 \mathcal V_E^\pi=
 \bigoplus_{\sum m r_m=E}
       \widehat{\bigotimes}_{\pi,m\geq1}
            \widehat{\operatorname{Sym}}_\pi^{r_m}H_m,
 \qquad \mathcal V^\pi=\bigoplus_E\mathcal V_E^\pi.
\end{equation}
Use the sum norm at each energy and the locally
convex direct-sum topology over energies.

\begin{proposition}\label{v4-prop:singular-projective-contractions}
Every internal and cross contraction is bounded
on the corresponding projective tensor sectors.
The canonical map \(\mathcal V_E^\pi\to\mathcal V_E\)
is injective with dense image. For real symmetric
quadratic coefficients,
\[
                       \|q_A\|_\pi=\sqrt2\|A\|_1,
\]
so this sector consists of symmetric trace-class
operators. The field operations, vacuum and
translation make \(\mathcal V^\pi\) a local
vertex algebra, and its map to \(\mathcal V\)
is a vertex homomorphism.
\end{proposition}

\begin{proof}
On a simple tensor, contraction multiplies the
remaining factors by a pairing bounded by
the two removed norms. Taking the infimum
over decompositions proves the bound. Iterate
it for a matching.

Finite species projections converge to the
identity in projective norm, first on simple
tensors and then by density and their uniform
bound. If a projective tensor maps to zero
in the Hilbert tensor product, each finite
species projection is zero. Its projective
limit is therefore zero. Finite species
tensors also prove density of the image.

A real symmetric spectral decomposition gives
the upper trace-norm bound. Pairing with the
bounded sign operator gives the reverse bound.
The Fock factor is \(\sqrt2\).
Each field coefficient at fixed energies has
finitely many bounded Wick patterns. The
finite-species identities extend by density,
with the uniform locality exponent already
proved. The two completions agree on their
common finite-species subspace, so their
inclusion preserves the coefficients.
\end{proof}

The unpaired mode substitution for a germ
\(f(0)=0\), \(f'(0)\ne0\), is
\begin{equation}\label{v4-eq:singular-coordinate-jets}
 T_f(m,j)=[z^{m-1}]f'(z)f(z)^{j-1}
             =\frac mj[z^m]f(z)^j,
 \qquad 1\leq j\leq m.
\end{equation}
The pair coefficient also has the expression
\begin{equation}\label{v4-eq:singular-coordinate-pairs}
 B_f(m,n)=mn[z^mu^n]\log\frac{f(z)-f(u)}{z-u}.
\end{equation}
Indeed two derivatives of the logarithm give
\(H_f\). Its constant is immaterial.

Write \(x_m(v)=a_{-m}(v)\mathbf1\).
Define \(R_f\mathbf1=\mathbf1\).

For a creation
monomial sum over all partial matchings of
its legs. An unpaired \(x_m(v)\) contributes
\(\sum_{j\leq m}T_f(m,j)x_j(v)\). A pair
\(x_m(v),x_n(w)\) contributes
\((v,w)B_f(m,n)\). 
Multiply within each
matching and sum. Energy at most \(E\)
uses only the \((E+1)\)-jet of \(f\).

\begin{theorem}\label{v4-thm:singular-projective-coordinates}
These formulas define bounded maps on every
bounded-energy part of \(\mathcal V^\pi\), with
\[
 R_fR_g=R_{f\circ g},\qquad
 R_f^{-1}=R_{f^{-1}}.
\]
For the recentered germ \(f_z(u)=f(z+u)-f(z)\),
they obey
\begin{align}
 R_fY(a,z)R_f^{-1}
   &=Y(R_{f_z}a,f(z)),\label{v4-eq:singular-coordinate-covariance}\\
 \partial_zR_{f_z}+f'(z)TR_{f_z}
   &=R_{f_z}T.\label{v4-eq:singular-coordinate-connection}
\end{align}
The identities are coefficientwise in the
coordinate variable.
\end{theorem}

\begin{proof}
Every matching is bounded by
Proposition~\ref{v4-prop:singular-projective-contractions}.
Unpaired jets compose by the chain rule.
In an iterated matching a pair is either
formed at the first change or formed after
that change. The kernel cocycle
\eqref{v4-eq:singular-kernel-cocycle} adds
these two contributions to the one-step
pair. Partitioning matchings in this way
proves composition at every tensor degree.
The identity and inverse germs give the
remaining group identities.

In finite species, the current transforms
as a one-form. Wick expansion of a composite
field adds exactly the difference of its
normal-ordering kernels. This is the same
matching rule, so it proves
\eqref{v4-eq:singular-coordinate-covariance}.
Continuity at fixed coefficients extends it
to projective tensors.

Applying it to the
vacuum gives
\(R_fe^{zT}=e^{f(z)T}R_{f_z}\).
Differentiation proves
\eqref{v4-eq:singular-coordinate-connection}.
These formal identities require no analytic
exponential on the algebraic energy sum.
\end{proof}

For \(n\geq-1\), the infinitesimal operators
on finite species are
\begin{equation}\label{v4-eq:singular-Witt-operators}
 L_n=\sum_{m>\max(n,0),i}
             m x_{m-n,i}\partial_{x_{m,i}}
    +\frac12\sum_{\substack{r+s=n\\r,s\geq1},i}
             rs\partial_{x_{r,i}}\partial_{x_{s,i}}.
\end{equation}
Here \(x_{m,i}=x_m(e_i)\). The sums on each
bounded-energy sector are continuous by the
contraction bounds. Direct commutation gives
\([L_m,L_n]=(m-n)L_{m+n}\), with zero bracket
when \(m=n=-1\).

The only possible central
cases have \(m+n=0\) and
\(m\in\{-1,0,1\}\), whose central coefficient
vanishes. Thus this is the indicated part
of the Witt action.

In particular
\[
 L_{-1}=T,
 \qquad L_0|_{\mathcal V_E^\pi}=E,
 \qquad L_2q_A=\operatorname{Tr}(A)\mathbf1.
\]
The operators \(L_{-1},L_0,L_1\) have no internal
contractions and also act on the Hilbert sectors.
Their existence does not require a state whose
field implements the full action.

\subsection{Descent to a curve}

On a smooth complex analytic curve, use holomorphic
sections which locally take values in a finite
sum of the Banach sectors \(\mathcal V_E^\pi\).
For \(w=f(z)\), glue their components by
\[
                        s_w(f(z))=R_{f_z}s_z(z).
\]
The finite jet matrices are holomorphic in
the root and locally bounded at each energy.
Their composition law proves the overlap
cocycle. Equation
\eqref{v4-eq:singular-coordinate-connection}
glues \(\nabla=d+T\,dz\). This connection
raises the local energy bound by one and
is flat on the curve.

Use completed projective external tensors,
finite pole orders and finite normal derivatives
for the local chiral operation. Its chart
formula is the principal part of \(h(u,z)Y(a,u)b\).
Tensoring the left module with \(\omega_X\)
gives the right differential module. A normal
class is still \(\delta\partial_u^r/r!\).
Covariance transports the expansion, density
and all normal jets together. Translation
and Jacobi therefore glue the analytic
topological chiral operation.

There is a separate algebraic version on a
smooth algebraic curve over \(\mathbb C\).
Regard \(\mathcal V^\pi\) as a vector space
and use algebraic external tensors. On a
state of energy at most \(E\), every word
in positive \(L_n\)'s with total index
greater than \(E\) vanishes. The finitely
many remaining words and the energy
projections span a finite-dimensional
stable subspace. Its coordinate action
is algebraic and factors through a finite
jet group.

Associate these subspaces to
the formal-coordinate torsor and take
their filtered union. The same identities
descend its connection and finite
principal-parts operation. The two
versions use different external tensors.

One may instead keep the Hilbert sectors
and choose a projective atlas, whose
transitions are fractional linear.
Their internal pair coefficients vanish,
so the same descent works with the
bounded Hilbert coefficient maps and
projective external tensors. The
projective atlas is part of this choice.
On \(\mathbb P^1\), inversion has the
recentered germ
\[
                    \frac1{p+u}-\frac1p=-\frac{u}{p(p+u)},
                    \qquad p\ne0,
\]
which is fractional linear. An affine
uniformizing atlas of an elliptic curve
gives another example. On a fixed such
atlas, \(\mathcal V^\pi\to\mathcal V\)
is dense, injective and compatible with
the operations, but is not a topological
equivalence in infinite species.

Finally let \(P\) be a projective connection,
with components
\(P_z=f'(z)^2P_w(f(z))+\{f,z\}\).
For symmetric trace-class \(A\),
\eqref{v4-eq:singular-coordinate-quadratic}
shows that
\[
 \left(W_A^{(z)}(z)
      -\frac{\operatorname{Tr}(A)}6P_z(z)\mathbf1\right)(dz)^2
\]
is coordinate independent. Changing \(P\)
changes this field by the corresponding
scalar quadratic differential. For a
positive non-trace-class \(A\), this
subtraction has no value in the stated
coefficient domain. It requires a trace,
an operation stronger than the cross
contractions which defined \(W_A\).

\section{A completion retaining the contractions}
\label{v4-sec:singular-contraction-completion}

The sequence \(q_{P_N/N}\) had Hilbert limit
zero and trace one. To retain both limits,
complete the finite-species states in the
norm of their state tensor together with
all their internal contractions. The new
limit will be a central state, distinct
from zero and from the vacuum.

Introduce symmetric variables
\[
 c_{mn}=c_{nm},\qquad \deg c_{mn}=m+n,
 \qquad \Gamma=\mathbb C[c_{mn}:1\leq m\leq n].
\]
Each homogeneous part \(\Gamma_d\) is finite
dimensional. Define
\begin{equation}\label{v4-eq:singular-contraction-space}
 \mathcal W_E=\bigoplus_{d=0}^E
                 \mathcal V_{E-d}\otimes\Gamma_d,
 \qquad \mathcal W=\bigoplus_{E\geq0}\mathcal W_E.
\end{equation}
Give central monomials an orthonormal basis
and use the finite Hilbert sum at each energy.
Any choice of norms on these finite-dimensional
central factors gives the same topology.

On finite-species oscillator polynomials put
\begin{equation}\label{v4-eq:singular-contraction-embedding}
 \Delta_c=\frac12\sum_{m,n\geq1}c_{mn}
                   \sum_i\partial_{x_m(e_i)}\partial_{x_n(e_i)},
 \qquad \sigma=e^{\Delta_c},\qquad j(a)=\sigma a.
\end{equation}
The exponential terminates on each polynomial.
Every derivative pair contracts two species
legs by their intrinsic pairing, so the
definition is independent of the orthonormal
basis. The norm \(\|j(a)\|_{\mathcal W_E}\)
is the norm of all its contraction coefficients.
Its constant coefficient in \(c\) is \(a\),
so it is indeed a norm.

\subsection{The completed field operations}

Define the central translation by
\begin{equation}\label{v4-eq:singular-central-translation}
 D_\Gamma c_{mn}=m c_{m+1,n}+n c_{m,n+1}.
\end{equation}
The translation of \(\mathcal W\) is
\(\mathsf T=T_x+D_\Gamma\), where
\(T_xx_m(v)=m x_{m+1}(v)\). For two modes set
\begin{align}
 P_m(z;v)&=\sum_{r\geq0}\binom{m+r-1}{r}
                       x_{m+r}(v)z^r,
                    \label{v4-eq:singular-contraction-Taylor}\\
 \kappa_{mn}(z)&=
   \frac{(-1)^{m-1}(m+n-1)!}{(m-1)!(n-1)!}\,z^{-m-n},
                    \label{v4-eq:singular-contraction-pole}\\
 G_{mn}(z)&=\sum_{r\geq0}\binom{m+r-1}{r}
                            c_{m+r,n}z^r.
                    \label{v4-eq:singular-contraction-regular}
\end{align}
The scalar \(\kappa_{mn}\) is the level-one
contraction already computed. The new term
\(G_{mn}\) retains its regular contraction data.

Define \(Y_{\mathcal W}(a,z)b\) by matchings
between the oscillator legs of its two
polynomial inputs. A pair of modes \(m,n\)
and species \(v,w\) contributes
\begin{equation}\label{v4-eq:singular-contraction-pair}
                    (v,w)\bigl(\kappa_{mn}(z)+G_{mn}(z)\bigr).
\end{equation}
An unpaired left leg contributes \(P_m(z;v)\).
An unpaired right leg stays at the root.

Translate the left central coefficient by
\(e^{zD_\Gamma}\) and multiply it by the
right one. The product of these terms,
summed over partial matchings, is the field.
Only contractions between separate Hilbert
tensors occur in this rule.

\begin{proposition}\label{v4-prop:singular-contraction-vertex}
The coefficient maps make \(\mathcal W\) a
local vertex algebra with vacuum \(1\) and
translation \(\mathsf T\). Its fixed-energy
mode maps are bounded bilinear, with locality
exponent \(E+F\) on input energies \(E,F\).
Both the finite-species map \(j\) and the
projection
\[
                  \pi:\mathcal W\longrightarrow\mathcal V,
                  \qquad c_{mn}\longmapsto0,
\]
are vertex homomorphisms on their stated domains.
\end{proposition}

\begin{proof}
On finite species start with the tensor product
of the Heisenberg vertex algebra and the
commutative differential algebra
\((\Gamma,D_\Gamma)\). Its field on central
coefficients is \((e^{zD_\Gamma}\gamma)\gamma'\).
Conjugate its field operation by \(\sigma\),
whose polynomial inverse is \(e^{-\Delta_c}\).
The commutator
\([T_x,\partial_{x_m}]=-(m-1)\partial_{x_{m-1}}\)
shows, after reindexing, that
\([T_x,\Delta_c]=-[D_\Gamma,\Delta_c]\).
Thus \(\sigma\) commutes with \(\mathsf T\).

In the conjugated product, the internal
contractions of \(\sigma^{-1}\) on each input
cancel those of \(\sigma\) on that same
input inside the output. The remaining
pairs connect the two inputs. A translated
left leg gives the new term \(G_{mn}\),
and the old singular pair gives
\(\kappa_{mn}\). This proves precisely
\eqref{v4-eq:singular-contraction-pair}.

At fixed energies and output coefficient,
there are finitely many pairings, Taylor
terms and central monomials. The cross
contraction bound proves their continuity.
The singular order is at most \(E+F\),
independently of species rank. The
finite-species identities therefore extend
by density as in
Theorem~\ref{v4-thm:singular-hilbert-vertex}.
Conjugation proves the assertion about \(j\).
Evaluation at \(c=0\) proves the assertion
about \(\pi\).
\end{proof}

\begin{theorem}\label{v4-thm:singular-contraction-completion}
The energywise completion of the finite-species
states in the norm
\(a\mapsto\|j(a)\|_{\mathcal W_E}\) is
exactly \(\mathcal W\).
The projection \(\pi\) is surjective and
has the bounded linear section \(v\mapsto v\otimes1\),
which is not a vertex homomorphism.
The map \(j\) extends to a continuous,
dense, proper vertex embedding
\(\mathcal V^\pi\hookrightarrow\mathcal W\).
\end{theorem}

\begin{proof}
For the states \(u_N^{mn}\) used in the
coordinate obstruction,
\[
 j(u_N^{mn})=u_N^{mn}+c_{mn},\qquad
 \|u_N^{mn}\|^2=
 \begin{cases}
   2m^2/N,&m=n,\\
   mn/N,&m\ne n.
 \end{cases}
\]
Their limits put every \(c_{mn}\) in the
closure of the image of \(j\). That closure
is stable under the continuous field
coefficients. On central states the
\((-1)\)-product is polynomial multiplication,
so it contains all central monomials.

For an oscillator polynomial \(a\), the
difference \(j(a)-a\) is a sum of central
monomials times polynomials with fewer
oscillator legs. Induction on the number
of legs therefore puts \(a\), and then
every \(ac^\nu\), in the closure. Such
polynomials are dense in every finite
Hilbert sum \(\mathcal W_E\). This proves
the completion assertion.

The linear section is not a vertex map.
For a unit species, the completed
current \((-1)\)-product is \(x_1^2+c_{11}\).
The section sends its old product to \(x_1^2\),
omitting the regular contraction.

Every internal contraction is bounded on
projective tensors, so \(j\) extends
continuously to \(\mathcal V^\pi\).
Its composite with \(\pi\) is the
injective projective-to-Hilbert map.
It is therefore injective. It has dense
image because it contains the finite-species
image of \(j\). Its composite with \(\pi\)
is injective, so the image contains no
nonzero \(c_{mn}\). The image is therefore proper.
\end{proof}

In the quadratic mode-one sector, this
embedding is the dense graph
\[
 \mathcal S_1^{\mathrm{sym}}(H)
   \longrightarrow\mathcal S_2^{\mathrm{sym}}(H)\oplus\mathbb Cc_{11},
 \qquad A\longmapsto(A,\operatorname{Tr}A).
\]
Here \(\mathcal S_p\) denotes the Schatten
class with summable \(p\)-th powers of
singular values. The completion has an
independent second coordinate. It has
not defined a trace on every old
Hilbert--Schmidt operator.

\subsection{The central algebra and ordinary actions}

\begin{theorem}\label{v4-thm:singular-contraction-centre}
The ordinary vertex center of \(\mathcal W\)
is \(\Gamma\). With its translation it is
the free commutative differential algebra
on one generator of each positive even energy:
\begin{equation}\label{v4-eq:singular-contraction-centre}
 \Gamma=\mathbb C[D_\Gamma^r\chi_k:k\geq1,\ r\geq0],
 \qquad \chi_k=c_{kk},\qquad \deg\chi_k=2k.
\end{equation}
Here a differential algebra means an
algebra with a derivation. \(D_\Gamma\)
is not a square-zero differential.
\end{theorem}

\begin{proof}
The regular term \(G_{mn}\) does not change
the positive current modes. They remain
\(a_n(v)=n\partial_{x_n(v)}\), \(n>0\).
Their common kernel on Hilbert oscillator
tensors is the vacuum. Since the central
monomial basis at each energy is finite,
the same argument coefficient by coefficient
puts every central state in \(\Gamma\).
Conversely its elements are central in
the transported finite-species algebra.
Continuity extends centrality to all of
\(\mathcal W\).

Identify the linear span of the \(c_{mn}\)
with \(\mathbb C[u,v]^{\mathfrak S_2}\) by
\[
 c_{mn}\longmapsto
 \frac{u^{m-1}v^{n-1}+u^{n-1}v^{m-1}}
      {2(m-1)!(n-1)!}.
\]
Under this map \(D_\Gamma\) is multiplication
by \(u+v\). The symmetric polynomial ring
\(\mathbb C[u+v,uv]\) is free over
\(\mathbb C[u+v]\) on the powers of \(uv\).
Up to nonzero factorials, these are the
images of \(c_{kk}\). Taking the symmetric
algebra proves both generation and
algebraic independence in
\eqref{v4-eq:singular-contraction-centre}.
\end{proof}

For instance,
\[
 c_{12}=\tfrac12D_\Gamma c_{11},\qquad
 c_{13}=\tfrac14(D_\Gamma^2c_{11}-2c_{22}).
\]
The additional even-energy generators are
needed to recover all two-mode contractions.

Let \(B=\bigoplus_{r\geq0}B_r\) be a
commutative energy-graded algebra with Banach
sectors, bounded multiplication and a bounded
degree-one derivation \(D_B\). Give it the
commutative vertex operation
\(Y_B(b,z)b'=(e^{zD_B}b)b'\).
In \(B\otimes\mathcal W\), complete each
external sector projectively and retain
the algebraic energy sum. An ordinary
action is an energy-preserving unital
vertex homomorphism
\(\rho:B\otimes\mathcal W\to\mathcal W\),
bounded on each sector and satisfying
\(\rho(1\otimes w)=w\).

\begin{proposition}\label{v4-prop:singular-central-actions}
These actions are naturally in bijection
with bounded graded algebra maps
\(\phi:B\to\Gamma\) commuting with the
two derivations. The corresponding action is
\[
                       \rho(b\otimes w)=\phi(b)_{(-1)}w.
\]
\end{proposition}

\begin{proof}
The two tensor factors have no singular
cross products. Since \(\rho\) fixes
\(\mathcal W\), its value on \(b\otimes1\)
is central and hence lies in \(\Gamma\).
Preservation of \((-1)\)-products and
translation gives the stated algebra map.
Conversely central multiplication gives
the displayed bounded vertex homomorphism.
The identity
\(b\otimes w=(b\otimes1)_{(-1)}(1\otimes w)\)
proves uniqueness. Both constructions
commute with maps of \(B\).
\end{proof}

This universal property fixes ordinary
continuous vertex actions and their
coefficient topology. A derived center
also requires its action spaces and
higher comparisons. A factorizing
transverse actor additionally requires
the specified coactions and collision
maps. Neither object is defined by
this ordinary-action property alone.

\subsection{Coordinates on the completed space}

Use the coordinate coefficients
\(T_f(m,j)\) and \(B_f(m,n)\) already
defined. On the polynomial coordinates
of \(\mathcal W\), set
\begin{align}
 R_fx_m(v)&=\sum_{j\leq m}T_f(m,j)x_j(v),
                    \label{v4-eq:singular-completed-current}\\
 R_fc_{mn}&=\sum_{j\leq m,\,k\leq n}
            T_f(m,j)T_f(n,k)c_{jk}+B_f(m,n).
                    \label{v4-eq:singular-completed-centre}
\end{align}
Extend by ordinary polynomial substitution.
At every energy bound, this uses only
finite-dimensional central maps and
finitely many bounded Hilbert mode maps.

\begin{theorem}\label{v4-thm:singular-completed-coordinates}
These substitutions extend continuously
to \(\mathcal W\), preserve bounded energy,
and satisfy
\[
 R_fj=jR_f^{\mathrm{fin}},\qquad
 R_fR_g=R_{f\circ g},\qquad R_f^{-1}=R_{f^{-1}}.
\]
Here \(R_f^{\mathrm{fin}}\) is the finite-species
matching formula. They also satisfy
\begin{align}
 R_fY_{\mathcal W}(a,z)R_f^{-1}
   &=Y_{\mathcal W}(R_{f_z}a,f(z)),
                   \label{v4-eq:singular-completed-covariance}\\
 \partial_zR_{f_z}+f'(z)\mathsf T R_{f_z}
   &=R_{f_z}\mathsf T.
                   \label{v4-eq:singular-completed-connection}
\end{align}
Thus locally bounded-energy holomorphic
sections glue on a smooth complex curve,
with connection \(d+\mathsf T\,dz\) and
their continuous field coefficients.
\end{theorem}

\begin{proof}
Expand \(j(a)\) into its self-pair matchings.
For each pair, \eqref{v4-eq:singular-completed-centre}
chooses either the transported central
coefficient or the scalar change in the
vacuum kernel. Summing these choices
is exactly the finite-species coordinate
matching followed by \(j\). This proves
the first identity. The kernel and
unpaired-jet cocycles prove composition,
and the inverse germ proves inversion.

The field identity holds on the dense
finite-species image by its established
covariance. All coefficients are bounded
at fixed energies, so it holds on
\(\mathcal W\). Applying it to the
vacuum and differentiating gives the
connection identity. Holomorphic dependence
of the finite jets, their composition
law and that identity glue the claimed
sections and flat connection.
\end{proof}

For the generating series
\(\mathcal C(z,u)=\sum_{m,n\geq1}c_{mn}z^{m-1}u^{n-1}\),
the central law reads
\begin{equation}\label{v4-eq:singular-central-bidifferential}
 R_f\mathcal C(z,u)=
       f'(z)f'(u)\mathcal C(f(z),f(u))+H_f(z,u).
\end{equation}
In particular
\[
 R_fc_{11}=f'(0)^2c_{11}+\tfrac16\{f,0\},
 \qquad L_2c_{11}=1.
\]
The infinitesimal operators for \(n\geq-1\)
are derivations of these polynomial coordinates:
\[
 \begin{aligned}
 L_nx_m(v)&=\mathbf1_{m>n}\,m x_{m-n}(v),\\
 L_nc_{pq}&=\mathbf1_{p>n}\,p c_{p-n,q}
           +\mathbf1_{q>n}\,q c_{p,q-n}
           +pq\,\mathbf1_{p+q=n}.
 \end{aligned}
\]
Excluded terms are omitted. For \(n=-1\)
this is \(\mathsf T\). Differentiating
the coordinate law gives the same
positive-half Witt brackets.

\begin{proposition}\label{v4-prop:singular-vacuum-radical}
The old Fock form has a unique continuous
extension along \(j\), namely
\[
             \langle a,b\rangle_0
                  =\langle\pi a,\pi b\rangle_{\mathrm{Fock}}.
\]
Its radical is \(\ker\pi\), which is not
preserved by all coordinate changes.
No conformal vector in \(\mathcal W\)
implements the displayed operator \(L_2\).
\end{proposition}

\begin{proof}
The form agrees with the old one on
\(j(\mathcal V_{\mathrm{fin}})\). Density
proves uniqueness on bounded-energy
sectors. Its radical is the vertex ideal
generated by the positive-degree central
variables. But \(L_2c_{11}=1\) sends
one such variable to the vacuum.
Every singular mode of every field
annihilates a central state. A conformal
vector implementing \(L_2\) would
therefore give \(L_2c_{11}=0\), a contradiction.
\end{proof}

In particular no positive-definite Hilbert
form on this completion extends that
specified vacuum form. A separately chosen
state with different values on \(\Gamma\)
is a different question.

\subsection{A chosen bidifferential and its quotients}

Choose a symmetric local bidifferential
\[
 \mathcal B(s,t)=\bigl((s-t)^{-2}+b(s,t)\bigr)ds\,dt,
\]
with \(b\) regular near the diagonal.
Its divided jets at the moving root \(p\) are
\[
 b_{mn}(p)=
 \left.\frac{\partial_s^{m-1}\partial_t^{n-1}b(s,t)}
                {(m-1)!(n-1)!}\right|_{s=t=p}.
\]
Take formal sections \(\mathcal W_E[[p]]\),
or holomorphic sections locally in finitely
many energy sectors. Scalar functions
carry the commutative vertex operation
with translation \(\partial_p\).

\begin{proposition}\label{v4-prop:singular-bidifferential-quotient}
The relations \(c_{mn}=b_{mn}(p)\) generate
a vertex ideal stable under
\(\mathsf T+\partial_p\). The quotient
has the coefficient space of locally
bounded-energy \(\mathcal V\)-valued
sections, with the regular contractions
\(G_{mn}\) evaluated at these jets.
The quotients glue wherever \(\mathcal B\) glues.
\end{proposition}

\begin{proof}
Differentiating both variables at the root gives
\[
                  \partial_pb_{mn}
                       =m b_{m+1,n}+n b_{m,n+1}.
\]
Translation therefore takes each relation
to a sum of such relations. Their
coefficients are central, so the ideal
is a vertex ideal. Evaluation in the
finite central polynomial at each energy
gives the quotient. Its kernel is this
ideal: subtract the evaluated value of
a monomial one factor at a time.
The transformation law of the regular
part of \(\mathcal B\) is
\eqref{v4-eq:singular-central-bidifferential},
which proves compatibility on overlaps.
\end{proof}

On \(\mathbb P^1\), the bidifferential
\(ds\,dt/(s-t)^2\) has \(b=0\) in
projective coordinates and recovers
the Hilbert construction with that
atlas. A choice on another curve is
additional geometric data.

Two choices \(b,b'\) need not have a
continuous current-preserving comparison
on the unchanged Hilbert sectors. Such
a comparison must send a finite-rank
quadratic state by
\begin{equation}\label{v4-eq:singular-background-change}
                q_A\longmapsto
                    q_A+\operatorname{Tr}(A)(b'_{11}-b_{11}),
\end{equation}
as follows by comparing the current
\((-1)\)-products. At a root where
the jet difference is nonzero, the
sequence \(q_{P_N/N}\) proves
nonclosability. Higher mode pairs
test the other jet differences.

Both quotient families exist. Their
specified comparison requires a
stronger tensor domain. On
\(\mathcal V^\pi\), all the needed
contractions are bounded and the
comparison does extend.

\section{Gaussian preparations and their observables}
\label{v4-sec:singular-gaussian}

A vacuum specifies the field algebra. A preparation
instead specifies a normalized vector and the
observables on which its expectations exist.
We construct such vectors in the first oscillator
of \(\widehat{\mathcal V}\), and then compute
their scalar field matrix elements.

\subsection{The first oscillator and Gaussian multiplication}

Choose independent standard real Gaussians
\((\xi_i)\) in a real orthonormal basis \((e_i)\) of \(H\).
The \(L^2\) sum
\(\xi(h)=\sum_i(h,e_i)\xi_i\) has covariance
\(\mathbb E[\xi(h)\xi(k)]=(h,k)\).
Let \(\mathcal F_\xi\) be the completed sigma field
they generate. For a symmetric tensor \(a\) of order
\(p\), define its multiple Gaussian integral \(I_p(a)\)
by polarization and the generating identity
\[
 e^{t\xi(h)-t^2\|h\|^2/2}
       =\sum_{p\geq0}\frac{t^p}{p!}I_p(h^{\otimes p}).
\]
Multiplying two such exponentials and taking their
expectation gives
\begin{equation}\label{v4-eq:singular-chaos-isometry}
 \mathbb E[I_p(a)I_q(b)]
          =\delta_{pq}p!\langle a,b\rangle.
\end{equation}

This identity extends \(I_p\) to the symmetric Hilbert
tensor power. Complex linear extension gives the
corresponding complex \(L^2\) space.

Write \(a\widetilde\otimes_r b\) for contraction of
\(r\) legs followed by symmetrization. Comparing
coefficients in the same exponential identity gives
\begin{equation}\label{v4-eq:singular-chaos-product}
 I_p(a)I_q(b)
   =\sum_{r=0}^{\min(p,q)}
       r!\binom pr\binom qr
           I_{p+q-2r}(a\widetilde\otimes_r b).
\end{equation}
The coefficient chooses \(r\) legs on each side
and a bijection between them.
For finite tensors this proves the formula directly.

Under Hilbert tensor approximation, its right side
converges in \(L^2\), by
Lemma~\ref{v4-lem:singular-cross-contractions}.
Its left side converges in \(L^1\), by
Cauchy--Schwarz and \eqref{v4-eq:singular-chaos-isometry}.
The limits agree. Iterating the formula shows that
finite sums of these completed chaoses form an
algebra and belong to every finite \(L^p\) space.

Let \(\mathcal V^{(1)}\subset\mathcal V\) consist
of states using only mode-one creators, and let
\(\Pi_1\) be the projection onto this subspace.
The map
\[
 \Phi\bigl(a_{-1}(h_1)\cdots a_{-1}(h_p)\mathbf1\bigr)
       =I_p\bigl(\operatorname{Sym}(h_1\otimes\cdots\otimes h_p)\bigr)
\]
preserves its Fock norm. Its product comparison is
\begin{equation}\label{v4-eq:singular-gaussian-projection}
 \Phi\left(\left.\Pi_1Y(a,z)b\right|_{z=1}\right)
                         =\Phi(a)\Phi(b),
 \qquad a,b\in\mathcal V^{(1)}.
\end{equation}

Indeed each cross contraction contributes \(z^{-2}\).
The projection keeps only the constant mode-one
term of every uncontracted transported creator.
The result is the finite Laurent polynomial whose
coefficients occur in \eqref{v4-eq:singular-chaos-product}.
Thus projection precedes evaluation at \(z=1\).

For a unit current state \(j\), this reads
\[
 \Pi_1Y(j,z)j=j^2+z^{-2}\mathbf1,\qquad
 \Phi(j^2)=\xi^2-1.
\]
The contracted term restores ordinary Gaussian
multiplication. The resulting product on
\(\mathcal V^{(1)}\) is associative, since \(\Phi\)
transports it to that multiplication.
It differs from the vertex \((-1)\)-product.
In particular, \(\Pi_1\) is not a vertex homomorphism.

\begin{proposition}\label{v4-prop:singular-gaussian-unitary}
The map \(\Phi\) extends to a unitary
\[
 U:\mathcal F_1\xrightarrow{\sim}
       L^2(\Omega,\mathcal F_\xi,\mu),
\]
where \(\mathcal F_1\) is the Hilbert completion of
\(\mathcal V^{(1)}\). On finite-particle sectors,
\(a_{-1}(h)+a_1(h)\) becomes multiplication by \(\xi(h)\).
For a real symmetric Hilbert--Schmidt operator \(A\),
define
\[
 \mathsf Q_A=U^{-1}M_{I_2(A)}U,\qquad
 \operatorname{Dom}(\mathsf Q_A)
     =\{\psi:I_2(A)U\psi\in L^2\}.
\]
This is a closed self-adjoint operator. On each
finite-particle sector it equals
\begin{equation}\label{v4-eq:singular-gaussian-quadratic-operator}
 \sum_{i,j}A_{ij}\bigl(
      a_{-1}(e_i)a_{-1}(e_j)
      +2a_{-1}(e_i)a_1(e_j)+a_1(e_i)a_1(e_j)\bigr).
\end{equation}
\end{proposition}

\begin{proof}
The chaos isometry proves isometry. Hermite
polynomials are dense for finitely many Gaussian
coordinates. Conditional expectations onto
increasing finite coordinate sigma fields converge
in \(L^2(\mathcal F_\xi)\), proving surjectivity.
The one- and two-tensor product formulas give
the two operator identities in finite rank.

On the \(p\)-particle sector, pair creation and
pair annihilation have norm bounds
\[
 \sqrt{(p+1)(p+2)}\|A\|_2,\qquad
 \sqrt{p(p-1)}\|A\|_2.
\]
The middle term has norm at most \(2p\|A\|\).
Hence the finite-rank expressions converge
on every such sector. Multiplication by the
real measurable function \(I_2(A)\), on its
displayed maximal domain, is self-adjoint.
This specifies the closed operator and its domain.
\end{proof}

The sigma field in this unitary is essential.
An additional independent centered sign is
orthogonal to every Gaussian polynomial.
It does not belong to the range of \(U\)
in an enlarged probability space.

\subsection{Connected contractions of quadratic fields}

Put \(Q_A=I_2(A)\). For variables with all moments,
their joint cumulant is
\[
 \operatorname{cum}(Z_1,\ldots,Z_n)
  =\sum_{\pi}(-1)^{|\pi|-1}(|\pi|-1)!
            \prod_{B\in\pi}\mathbb E\prod_{i\in B}Z_i,
\]
where \(\pi\) ranges over partitions of
\(\{1,\ldots,n\}\). The same partition expression
defines a connected vacuum coefficient of fields.

\begin{proposition}\label{v4-prop:singular-quadratic-cycles}
For symmetric Hilbert--Schmidt \(A_1,\ldots,A_n\),
\(n\geq2\), one has
\begin{equation}\label{v4-eq:singular-quadratic-cumulant}
 \operatorname{cum}(Q_{A_1},\ldots,Q_{A_n})
   =2^{n-1}\sum_{\sigma\in S_{\{2,\ldots,n\}}}
       \operatorname{Tr}(A_1A_{\sigma(2)}\cdots A_{\sigma(n)}).
\end{equation}
Introduce symmetric independent variables \(k_{ij}=k_{ji}\).
For \(i_1=1\), \(i_j=\sigma(j)\) and \(i_{n+1}=i_1\),
the polynomial
\begin{equation}\label{v4-eq:singular-common-cycle-polynomial}
 2^{n-1}\sum_{\sigma\in S_{\{2,\ldots,n\}}}
      \operatorname{Tr}(A_{i_1}\cdots A_{i_n})
                         \prod_{j=1}^n k_{i_j i_{j+1}}
\end{equation}
gives this cumulant at \(k_{ij}=1\).
At \(k_{ij}=(z_i-z_j)^{-2}\), it gives the
connected vacuum coefficient of the quadratic
fields \(W_{A_i}(z_i)\).
\end{proposition}

\begin{proof}
The partition sum cancels disconnected Wick graphs.
Each remaining vertex has degree two, so the
connected graph is a cycle. For \(n\geq3\), an
unoriented cyclic order has \(2^n\) choices of
labeled legs. Its two orientations give the
coefficient \(2^{n-1}\) in the permutation sum.
Reversal preserves the trace because the
matrices are symmetric. For \(n=2\), there
are two pairings and a double edge.

The species sum around each cycle is the
displayed trace. Two Hilbert--Schmidt factors
have trace-class product. The other factors
are bounded. Thus each trace is continuous
under Hilbert--Schmidt approximation.
The Gaussian moments and the field coefficients
converge by their contraction bounds.
The two choices of edge kernel give the
two evaluations of the same polynomial.
\end{proof}

For three fields, the formula is
\[
 \langle W_A(z_1)W_B(z_2)W_C(z_3)\rangle_0
     =\frac{8\operatorname{Tr}(ABC)}
             {z_{12}^2z_{23}^2z_{31}^2},
 \qquad z_{ij}=z_i-z_j.
\]
Higher arity retains the order of noncommuting
coefficients inside each trace. The Gaussian
specialization is an evaluation of independent
edge variables. It does not require complex
positions with every pairwise contraction equal
to one.

\subsection{Normalized squeezed vectors}

Let \(K\) be a real symmetric Hilbert--Schmidt
operator with \(\|K\|<1\). Define the creation
series
\begin{equation}\label{v4-eq:singular-squeezed-vector}
 S_K=\exp\left(\frac12\sum_{i,j}K_{ij}
                a_{-1}(e_i)a_{-1}(e_j)\right)\mathbf1,
 \qquad
 \Omega_K=\det(1-K^2)^{1/4}S_K.
\end{equation}
The determinant is the Fredholm determinant:
\(K^2\) is trace class and \(1-K^2\) is
strictly positive. The exponential is first
interpreted in finite rank.

For real self-adjoint Hilbert--Schmidt \(A\)
with \(1+A\geq\epsilon1>0\), put
\begin{equation}\label{v4-eq:singular-gaussian-density}
 D_A=\det{}_2(1+A)^{1/2}e^{-I_2(A)/2},
 \qquad
 \det{}_2(1+A)=\prod_j(1+a_j)e^{-a_j}.
\end{equation}
Here \(a_j\) are its real eigenvalues.
The logarithms of the product factors have
a quadratic zero at zero, so they are summable.

\begin{theorem}\label{v4-thm:singular-squeezed-preparation}
The vectors in \eqref{v4-eq:singular-squeezed-vector}
exist in \(\mathcal F_1\), independently of the
finite-rank approximations, and
\[
       \|S_K\|^2=\det(1-K^2)^{-1/2},
       \qquad \|\Omega_K\|=1.
\]
The function \(D_A\) is a positive density
of integral one. With
\[
                         K=-A(2+A)^{-1},
\]
one has \(U\Omega_K=\sqrt{D_A}\).
For two such preparations, with kernels \(K,L\),
\begin{equation}\label{v4-eq:singular-squeezed-overlap}
 \langle\Omega_K,\Omega_L\rangle
   =\frac{\det(1-K^2)^{1/4}\det(1-L^2)^{1/4}}
           {\det(1-KL)^{1/2}}
   =\mathbb E\sqrt{D_AD_B}.
\end{equation}
The square root is positive and continued from
\(K=L=0\). No commutativity assumption is imposed.
\end{theorem}

\begin{proof}
For one oscillator the squared norm is
\[
 \sum_{m\geq0}\frac{(2m)!}{4^m(m!)^2}k^{2m}
                              =(1-k^2)^{-1/2}.
\]
Diagonalization proves the finite-rank norm
formula. The Hermite generating identity also gives
\[
 US_K=\det(1+K)^{-1/2}
          \exp\left(\frac12\xi^tK(1+K)^{-1}\xi\right).
\]

For the overlap of two such functions, the
Gaussian precision is
\[
 P=1-K(1+K)^{-1}-L(1+L)^{-1}.
\]
Each of the two subtracted self-adjoint operators
has upper bound strictly less than \(1/2\).
Thus \(P>0\), and direct multiplication gives
\[
                         (1+K)P(1+L)=1-KL.
\]
Gaussian integration proves
\(\langle S_K,S_L\rangle=\det(1-KL)^{-1/2}\).
This argument fixes both the order and the branch.

Finite-rank compressions converge in Hilbert--Schmidt
norm. Their squares and products converge in
trace norm, so the displayed determinants converge.
The overlap of two approximations to the same
kernel makes their squeezed vectors Cauchy.
It also proves independence of the approximation.
Their fixed-particle components are the
continuous tensor polynomials in \(K\) from
\eqref{v4-eq:singular-squeezed-vector}.

For finite-rank \(A_N\), the density is
\[
 D_{A_N}=\det(1+A_N)^{1/2}e^{-\xi^tA_N\xi/2},
 \qquad \mathbb E D_{A_N}=1.
\]
Choose \(p>1\) sufficiently close to one
that \(1+pA\) remains uniformly positive.
Compressions retain common spectral bounds, and
\[
 \mathbb E D_{A_N}^p
   =\det{}_2(1+A_N)^{p/2}\det{}_2(1+pA_N)^{-1/2}.
\]
The logarithm is uniformly bounded by a
constant times \(\|A_N\|_2^2\).
Indeed the scalar logarithmic ratio has
a quadratic zero at zero on the common
bounded spectral interval.

The chaos isometry gives \(I_2(A_N)\to I_2(A)\)
in \(L^2\). Regularized determinant convergence
follows from
\[
 \log(1+A)-A=-\int_0^1tA^2(1+tA)^{-1}\,dt.
\]
Squares converge in trace norm, and the
resolvents are uniformly bounded.

The densities therefore converge in probability.
The \(L^p\) bound makes them uniformly
integrable, so they converge in \(L^1\).
Their limit is \(D_A\) and has integral one.
Their positive square roots converge in \(L^2\).

On the bounded spectrum of \(A\), separated
from \(-1\), the function \(-a/(2+a)\)
has absolute value strictly below one.
It defines a Hilbert--Schmidt strict contraction.
For \(k=-a/(2+a)\), the one-coordinate identity is
\[
 (1-k^2)^{1/4}(1+k)^{-1/2}
       e^{k\xi^2/(2(1+k))}
                   =(1+a)^{1/4}e^{-a\xi^2/4}.
\]
The map \(A\mapsto-A(2+A)^{-1}\) is continuous
in Hilbert--Schmidt norm on this positive domain,
by the resolvent identity.
Passing to the two limits proves the square-root
formula. The overlap identity follows by unitarity.
\end{proof}

The prepared covariance of the quadratures is
\((1+A)^{-1}\). More generally, for a bounded
positive operator \(C\), put \(X(f)=\xi(C^{1/2}f)\).
Its covariance in this preparation is
\[
                     C^{1/2}(1+A)^{-1}C^{1/2}.
\]
Finite Gaussian characteristic functions give this
formula before taking limits. The \(L^1\) density
limit preserves those bounded characteristic functions.
The \(L^p\) estimates and Hölder's inequality also
preserve every fixed polynomial expectation.

\begin{proposition}\label{v4-prop:singular-non-HS-preparation}
Suppose \(A\geq0\) is compact but not Hilbert--Schmidt.
For \(K=-A(2+A)^{-1}\), the unnormalized squeezed
norm product diverges. The normalized finite
spectral preparations converge weakly to zero
and have no norm limit in \(\mathcal F_1\).
\end{proposition}

\begin{proof}
The eigenvalues \(a_j\to0\) give
\(|k_j|=a_j/(2+a_j)\), comparable to \(a_j\).
Thus \(\sum_jk_j^2=\infty\), and
\[
                  -\frac12\sum_j\log(1-k_j^2)=\infty.
\]
The normalized vacuum overlap of a tail,
\(\prod_j(1-k_j^2)^{1/4}\), tends to zero.

The overlap with any vector on finitely
many eigenmodes contains this tail factor.
Such vectors are dense, and all preparations
have norm one. Hence their weak limit is zero.
A norm limit would have norm one and the same
weak limit, a contradiction.
\end{proof}

\subsection{Prepared field matrix elements}

The vector \(\Omega_K\) usually has infinitely
many energies. We use it to define scalar
matrix elements, without defining \(Y(\Omega_K,z)\)
on all of \(\widehat{\mathcal V}\).

Let \(\mathsf N\) count particles in the first
oscillator. If \(s>0\) and \(e^{2s}\|K\|<1\), then
\begin{equation}\label{v4-eq:singular-prepared-number-bound}
 \|e^{s\mathsf N}\Omega_K\|^2
       =\det(1-K^2)^{1/2}
          \det(1-e^{4s}K^2)^{-1/2}<\infty.
\end{equation}
This follows from \(e^{s\mathsf N}S_K=S_{e^{2s}K}\).
The overlap proof also gives convergence of the
finite-rank preparations in this weighted norm.
It controls every particle moment.

Define
\[
 N_K=K^2(1-K^2)^{-1},\qquad
 M_K=K(1-K^2)^{-1}.
\]
The first operator is trace class. The second
is Hilbert--Schmidt. Writing
\(\langle O\rangle_K=\langle\Omega_K,O\Omega_K\rangle\),
the first-mode contractions are
\[
 \langle a_{-1}(v)a_1(u)\rangle_K=(v,N_Ku),
 \qquad
 \langle a_1(v)a_1(u)\rangle_K=(v,M_Ku).
\]

To see the full Wick rule, put
\(R=(1-K^2)^{-1/2}\) and
\[
             b(v)=a_1(Rv)-a_{-1}(KRv).
\]
These operators satisfy the canonical commutation
relations and annihilate \(\Omega_K\).
The identity
\(a_1(v)=b(Rv)+b(KRv)^*\) reduces its
finite-rank contractions to vacuum contractions.
The weighted bound passes each fixed polynomial
mode identity to the limit. Higher oscillator
modes remain in their vacuum.

\begin{theorem}\label{v4-thm:singular-prepared-fields}
For real species vectors \(v,u\), the scalar
current matrix element is
\begin{equation}\label{v4-eq:singular-prepared-current}
 \begin{aligned}
 \langle J_v(z)J_u(w)\rangle_K
   ={}&\frac{(v,u)}{(z-w)^2}
        +(v,N_Ku)(z^{-2}+w^{-2})\\
      &+(v,M_Ku)(1+z^{-2}w^{-2}).
 \end{aligned}
\end{equation}
For finitely many finite-energy states,
finite oscillator and species cutoffs of
their scalar matrix elements converge uniformly
on compact subsets of
\[
                         |z_1|>\cdots>|z_n|>0.
\]
The limits continue meromorphically.
Their possible poles lie on pairwise
diagonals and at the prepared punctures \(z_i=0\).
\end{theorem}

\begin{proof}
The vacuum current contraction is the convergent
series \(\sum_{m>0}m z^{-m-1}w^{m-1}=(z-w)^{-2}\).
The two creation--annihilation orderings in
mode one add the two \(N_K\) terms.
The two anomalous contractions add the
two \(M_K\) terms.
This proves \eqref{v4-eq:singular-prepared-current}.
The geometric series and every fixed derivative
converge uniformly on compact nested annuli.

Each finite-energy state has finitely many
occupancy types. Wick expansion gives finitely
many contraction patterns for those types.
Within one insertion, vacuum normal ordering
removes the identity self-contraction.
The remaining self kernels are derivatives
of the prepared-minus-vacuum kernel.
They are Hilbert--Schmidt in species.
Cross kernels are bounded species operators.

Absorb each bounded cross operator into
one tensor leg. Regard every self kernel
as a separate two-leg Hilbert tensor.
The resulting contraction graph has no
self edge. Merge its vertices successively,
contracting all edges between the two
groups at each step.
Lemma~\ref{v4-lem:singular-cross-contractions}
bounds the result by the tensor norms.

The bounds are uniform on compact nested annuli.
This proves the cutoff convergence.
The finitely many positional factors are
derivatives of the displayed rational kernel,
so they have precisely the stated possible poles.
\end{proof}

For a quadratic field, the self-contraction is
\begin{equation}\label{v4-eq:singular-prepared-quadratic}
 \langle W_A(z)\rangle_K
      =2z^{-2}\operatorname{Tr}(AN_K)
            +(1+z^{-4})\operatorname{Tr}(AM_K).
\end{equation}

Both products have traces. Let \(G_K(z,w)\)
be the species operator of the current kernel
in \eqref{v4-eq:singular-prepared-current}.
The connected quadratic \(n\)-point coefficient is
\[
 2^{n-1}\sum_{\sigma\in S_{\{2,\ldots,n\}}}
       \operatorname{Tr}
          \left(\prod_{j=1}^n
                  A_{i_j}G_K(z_{i_j},z_{i_{j+1}})\right),
\]
with the cyclic indices from
\eqref{v4-eq:singular-common-cycle-polynomial}.
The product retains its indicated order.
Every trace contains two Hilbert--Schmidt
\(A_i\) factors, which justify the same cycle proof.

For quadratures, the covariance instead becomes
\[
 1+2N_K+2M_K=(1+K)(1-K)^{-1}=(1+A)^{-1}.
\]
Thus the Gaussian and field comparisons concern
specified observables. The puncture terms also
show that the prepared state need not be
translation invariant.

\section{Transport of Gaussian preparations}
\label{v4-sec:singular-gaussian-transport}

Coordinate covariance on finite-energy states
does not itself transport an infinite-energy
preparation. The cubic coordinate flow isolates
the additional convergence problem:
\[
                  f_t(z)=\frac{z}{\sqrt{1-2tz^2}}.
\]
On finite-energy states its coordinate action
is \(e^{tL_2}\). On mode-one polynomials,
\[
                       L_2=\frac12\sum_i\partial_{x_i}^2.
\]
The sum is finite on each polynomial.

For a symmetric strict contraction \(K\), write
\[
                             K_t=K(1-tK)^{-1}.
\]
Every transport formula below uses a common
real interval about zero on which its resolvents
exist and all displayed response kernels remain
strict contractions. Scalar branches start
at one when \(t=0\).

\begin{theorem}\label{v4-thm:singular-cubic-transport}
If \(K\) is real symmetric and trace class,
finite-species cubic transport has the Fock limit
\begin{equation}\label{v4-eq:singular-cubic-transport}
 \lim_{N\to\infty}e^{tL_{2,N}}S_{P_NKP_N}
                  =\det(1-tK)^{-1/2}S_{K_t}.
\end{equation}
Here \(P_N\to1\) strongly are finite-rank
orthogonal projections, and \(L_{2,N}\)
contracts only species in \(P_NH\).
If \(K\geq0\) is Hilbert--Schmidt but not
trace class, the vacuum coefficients of
these outputs diverge for sufficiently small \(t>0\).
\end{theorem}

\begin{proof}
In finite rank, insert
\(F_t\exp(x^tK_tx/2)\) in the heat equation.
Comparison of its quadratic and constant
terms gives
\[
           \dot K_t=K_t^2,\qquad
           \frac{d}{dt}\log F_t=\frac12\operatorname{Tr}K_t.
\]
The initial conditions yield the displayed
resolvent and ordinary determinant.
Their Gaussian series converges in Fock norm
on the stipulated interval.

Trace-class compressions converge in trace norm.
The resolvent identity gives convergence of
their responses in that norm.
Determinant continuity and the squeezed
overlap formula prove \eqref{v4-eq:singular-cubic-transport}.

For positive \(K\) and small positive \(t\),
\[
 \log\det(1-tP_NKP_N)^{-1/2}
                  \geq\frac t2\operatorname{Tr}(P_NKP_N).
\]
The traces tend to infinity if \(K\) is
not trace class. Indeed strong convergence
first captures any fixed finite spectral
subspace. Those subspaces have unbounded traces.
The determinant is the vacuum coefficient,
so the proposed Hilbert limit cannot exist.
\end{proof}

For negative kernels the same normalization
can have a zero limit instead of diverging.
If \(K\leq0\) is Hilbert--Schmidt but not
trace class, then for small \(t>0\),
\[
                         e^{tL_{2,N}}S_{P_NKP_N}\longrightarrow0.
\]

To prove this, put \(K_N=P_NKP_N\) and
\(q=\|K\|\). The response kernels converge
in Hilbert--Schmidt norm and have a common
strict-contraction bound. Their squeezed norms
are consequently bounded. Meanwhile
\[
 \log\det(1-tK_N)^{-1/2}
       \leq-\frac{t}{2(1+tq)}\operatorname{Tr}|K_N|
                                                \longrightarrow-\infty,
\]
using \(\log(1+x)\geq x/(1+x)\).
For small \(t<0\), the reverse sign gives
diverging vacuum coefficients.

Normalized inputs \(\Omega_{K_N}\) have the
same zero positive-time limit, since their
normalizing constants have a positive limit.
Normalizing each output separately instead
gives \(\Omega_{(K_N)_t}\to\Omega_{K_t}\).
This last operation changes the scalar
normalization. It does not make the
zero limit invertible.

These limits concern the specified Gaussian
families. They do not construct a closable
generator on full Hilbert Fock space.
The finite-energy sequence in
Theorem~\ref{v4-thm:singular-coordinate-obstruction}
still obstructs that extension.

\subsection{A regularized scalar and its composition defect}

For Hilbert--Schmidt \(K\), define
\[
 \widehat F_t(K)=\det{}_2(1-tK)^{-1/2},
 \qquad
 \log\widehat F_t(K)
       =\frac12\sum_{r\geq2}\frac{t^r}{r}\operatorname{Tr}(K^r).
\]
The series converges when \(\|tK\|<1\).
It omits the possibly undefined term
\(t\operatorname{Tr}(K)/2\).
That omission changes composition.

\begin{proposition}\label{v4-prop:singular-determinant-defect}
On the common domain for \(s,t,s+t\),
\begin{equation}\label{v4-eq:singular-determinant-defect}
 \widehat F_t(K)\widehat F_s(K_t)
  =\widehat F_{t+s}(K)
     \exp\left[-\frac{st}{2}
                 \operatorname{Tr}\bigl(K^2(1-tK)^{-1}\bigr)\right].
\end{equation}
The multiplier depends on \(K\).
\end{proposition}

\begin{proof}
The resolvent identities give
\[
 (1-sK_t)(1-tK)=1-(s+t)K,\qquad
 K_t-K=tK^2(1-tK)^{-1}.
\]
In finite rank, ordinary determinants multiply.
Their regularizing exponentials leave
\(-s\operatorname{Tr}(K_t-K)/2\) in the logarithm.
This is the displayed exponent.

Hilbert--Schmidt approximation gives trace-norm
convergence of the squared kernels and
uniformly bounded resolvents.
The logarithmic series converge locally uniformly,
so the same identity holds in the stated domain.
At \(K=0\) the multiplier is one.
For nonzero positive rank-one \(K\) and
positive \(s,t\), it is not one.
\end{proof}

Consequently this multiplier cannot be removed
on every squeezed line by a normalization
depending only on the coordinate parameters.
A map assigned to each such line also
does not define a linear map on Fock space.

A chosen reference preparation gives a
coherent comparison of lines along its orbit.
Fix \(K_0\), put \(K_t=K_0(1-tK_0)^{-1}\), and set
\[
 g(t)=\widehat F_t(K_0),\qquad
 \mathcal R(s,t)=g(t)/g(s).
\]
The difference \(K_s-K_0\) is trace class, and
\eqref{v4-eq:singular-determinant-defect} gives
\begin{equation}\label{v4-eq:singular-anchored-transport}
 \mathcal R(s,t)
   =\widehat F_{t-s}(K_s)
        \exp\left[\frac{t-s}{2}\operatorname{Tr}(K_s-K_0)\right].
\end{equation}
By its ratio definition,
\(\mathcal R(t,u)\mathcal R(s,t)=\mathcal R(s,u)\).
This construction retains the reference
preparation and its trace correction.

For a general coordinate germ, only two
coefficients act on mode-one polynomials:
\[
                   \lambda=f'(0),\qquad b=B_f(1,1)=\{f,0\}/6.
\]
Assume these numbers are real and \(\lambda\ne0\).
Let \(D_\lambda\) multiply a creation monomial
of length \(p\) by \(\lambda^p\).
Each paired leg contributes \(b\), and
each unpaired leg contributes \(\lambda\).
Thus \(R_f=D_\lambda e^{bL_2}\) on those polynomials.

For trace-class \(K\), the corresponding cutoff limit is
\begin{equation}\label{v4-eq:singular-first-mode-transport}
 R_fS_K=\det(1-bK)^{-1/2}
                   S_{\lambda^2K(1-bK)^{-1}}.
\end{equation}
It has the same common-domain qualifications
as the cubic formula.
When \(b=0\), the formula holds for
Hilbert--Schmidt \(K\), provided
\(\|\lambda^2K\|<1\).

Kernel composition gives
\[
 b_{f\circ g}=b_g+\lambda_g^2b_f,\qquad
 \lambda_{f\circ g}=\lambda_f\lambda_g.
\]
These data determine the first-oscillator
transport. Higher oscillator modes still
detect the other coefficients \(B_f(m,n)\).

\subsection{Retaining a trace with the preparation}

The contraction completion permits a more
general transport datum. Retain modes
\(1,\ldots,d\), with \(d<\infty\), and put
\[
 y_m(v)=x_m(v)/\sqrt m,\qquad
 \mathcal H=\mathbb R^d\otimes H,\qquad
 C_{mn}=c_{mn}/\sqrt{mn}.
\]
The central matrix \(C\) acts identically
on species. For trace-class block operators
on \(\mathcal H\), define the partial trace
\[
                   \operatorname{tr}_H(A)_{mn}
                             =\operatorname{Tr}_H(A_{mn}).
\]
It is bounded from the trace class to
the finite matrix space.
No such trace is assigned to an arbitrary
Hilbert--Schmidt block operator.

Fix a real symmetric Hilbert--Schmidt strict
contraction \(K\) on \(\mathcal H\).
Choose independently a real symmetric matrix
\(\tau\) of size \(d\).
If \(K\) is trace class, its canonical
choice is \(\tau=\operatorname{tr}_H(K)\).

\begin{lemma}\label{v4-lem:singular-trace-approximation}
Every pair \((K,\tau)\) has finite-rank
symmetric approximations \(K_N\) with
\[
          \|K_N-K\|_2\longrightarrow0,\qquad
          \operatorname{tr}_H(K_N)=\tau.
\]
The approximations are uniformly strict
contractions for sufficiently large \(N\).
If \(K\geq0\) is not trace class,
approximations with bounded prescribed
traces cannot all be positive.
\end{lemma}

\begin{proof}
Compress \(K\) to the first \(N\) species
in each mode, obtaining \(K_N^0\).
Put \(\delta_N=\tau-\operatorname{tr}_H(K_N^0)\).
Choose an orthogonal reservoir of \(M_N\)
further species and set
\[
 K_N=K_N^0+\delta_N\otimes P_{\mathrm{res}}/M_N,
 \qquad M_N>N^2\|\delta_N\|_F^2+N+1.
\]
The added term has partial trace \(\delta_N\)
and Hilbert--Schmidt norm
\(\|\delta_N\|_F/\sqrt{M_N}<1/N\).
Hence \(K_N\to K\) in Hilbert--Schmidt norm,
and therefore in operator norm.
The strict-contraction bound follows.

For positive approximations and any fixed
finite-rank projection \(Q\),
\[
 \operatorname{Tr}(QKQ)
    =\lim_N\operatorname{Tr}(QK_NQ)
               \leq\liminf_N\operatorname{Tr}K_N.
\]
The supremum of the left side is infinite
for a positive non-trace-class \(K\).
Thus bounded prescribed traces require
signed reservoir corrections in that case.
\end{proof}

In the product of energy sectors define
\begin{equation}\label{v4-eq:singular-gaussian-lift}
 \begin{aligned}
 \mathscr S_{K,\tau}(y,C)
  ={}&\exp\left[
       \frac12\operatorname{tr}_d(C\tau)
       +\frac12\sum_{r\geq2}\frac{\operatorname{Tr}((CK)^r)}r
             \right]\\
    &\quad\cdot
       \exp\left[\frac12y^tK(1-CK)^{-1}y\right].
 \end{aligned}
\end{equation}
The inverse is expanded in the
positive-energy central variables.
The traces for \(r\geq2\) exist
because each contains two Hilbert--Schmidt factors.

\begin{proposition}\label{v4-prop:singular-gaussian-lift}
The formula defines an element of
\(\prod_E\mathcal W_E\).
For every sequence from
Lemma~\ref{v4-lem:singular-trace-approximation},
it is the coefficientwise limit of \(jS_{K_N}\).
Under the sectorwise extension of \(\pi\),
it gives the Hilbert squeezed vector \(S_K\).
\end{proposition}

\begin{proof}
In finite rank the Gaussian heat equation gives
\[
 e^{\Delta_c}e^{y^tKy/2}
       =\det(1-CK)^{-1/2}
                     e^{y^tK(1-CK)^{-1}y/2}.
\]
Separating the determinant's linear term
gives \eqref{v4-eq:singular-gaussian-lift}.
Closed cycles of length at least two
have trace-class products.
Open paths give Hilbert--Schmidt kernels,
whose products give the required Hilbert tensors.

At fixed energy only finitely many paths,
cycles and oscillator degrees occur.
Their coefficients converge in trace norm
or Hilbert tensor norm.
The one-edge cycles converge by the
prescribed partial trace \(\tau\).
This proves the limit and its independence
of the approximating sequence.
At \(C=0\) the formula is \(S_K\).
Theorem~\ref{v4-thm:singular-squeezed-preparation}
applies to the one-particle space \(\mathcal H\).
\end{proof}

\subsection{Response, trace and scalar under composition}

A finite-mode Gaussian transformation consists
of \(T\in\mathrm{GL}_d(\mathbb R)\) and
\(B\in\operatorname{Sym}_d(\mathbb R)\).
On polynomials it substitutes \(y\mapsto Ty\)
after contracting pairs with \(B\).
On the completed polynomial coordinates
this is
\[
                   y\longmapsto Ty,\qquad
                   C\longmapsto TCT^t+B.
\]
A coordinate germ gives such a pair by
\[
 T_{ij}=\sqrt{j/i}\,T_f(i,j),\qquad
 B_{ij}=B_f(i,j)/\sqrt{ij}.
\]
An arbitrary matrix pair need not arise
from a curve coordinate.

Write \(\widetilde T=T\otimes1_H\) and
\(\widetilde B=B\otimes1_H\). Define
\begin{align}
 K^B&=K(1-\widetilde BK)^{-1},
                 \label{v4-eq:singular-trace-resolvent}\\
 K'&=\widetilde T^tK^B\widetilde T,
                 \label{v4-eq:singular-trace-response}\\
 \tau'&=T^t\tau T+
       \operatorname{tr}_H\bigl(\widetilde T^t(K^B-K)\widetilde T\bigr),
                 \label{v4-eq:singular-trace-update}\\
 F_{T,B}(K,\tau)&=
       e^{\operatorname{tr}_d(B\tau)/2}
               \det{}_2(1-\widetilde BK)^{-1/2}.
                 \label{v4-eq:singular-trace-scalar}
\end{align}
The resolvent identity
\[
 K^B=(1-K\widetilde B)^{-1}K
\]
shows that the response is symmetric.
The partial trace in
\eqref{v4-eq:singular-trace-update} is legitimate:
\begin{equation}\label{v4-eq:singular-trace-difference}
 K^B-K=K(1-\widetilde BK)^{-1}\widetilde BK
                                      \in\mathcal S_1(\mathcal H).
\end{equation}
Its two outside kernels are Hilbert--Schmidt.

Use a common domain where
\(\|\widetilde BK\|<1\) and every response
kernel is a strict contraction.
These conditions must hold for each
intermediate and composite transformation.
The logarithmic determinant series beginning
at degree two fixes its branch at the identity.

\begin{theorem}\label{v4-thm:singular-trace-transport}
On this domain, the finite-rank transformations
have the Gaussian cutoff limit
\[
       \mathscr S_{K,\tau}
             \longmapsto F_{T,B}(K,\tau)\mathscr S_{K',\tau'}.
\]
If \((T_2,B_2)\) is applied first, followed
by \((T_1,B_1)\), their composite data are
\begin{equation}\label{v4-eq:singular-gaussian-pair-law}
                   T=T_2T_1,\qquad B=B_2+T_2B_1T_2^t.
\end{equation}
The one-step and two-step responses and
trace matrices agree, and
\begin{equation}\label{v4-eq:singular-gaussian-scalar-law}
 F_{T_2,B_2}(K,\tau)
     F_{T_1,B_1}(K_2,\tau_2)=F_{T,B}(K,\tau).
\end{equation}
These identities hold under every finite
regrouping of admitted transformations.
No commutativity assumption is needed.
\end{theorem}

\begin{proof}
Finite-dimensional Gaussian heat integration
followed by linear substitution gives
\eqref{v4-eq:singular-trace-response} and the
scalar \(\det(1-\widetilde BK)^{-1/2}\).
Separating its linear trace term gives
\eqref{v4-eq:singular-trace-scalar} when
\(\tau=\operatorname{tr}_H K\).
Taking the partial trace of the response
gives \eqref{v4-eq:singular-trace-update}.

Apply the approximations of
Lemma~\ref{v4-lem:singular-trace-approximation}.
Resolvents converge in operator norm,
responses in Hilbert--Schmidt norm, and
\eqref{v4-eq:singular-trace-difference}
in trace norm.
For \(r\geq2\), the logarithmic series obeys
\[
 |\operatorname{Tr}((\widetilde BK)^r)|
       \leq\|\widetilde BK\|_2^2\|\widetilde BK\|^{r-2}.
\]
It converges locally uniformly.
Hence the scalar and trace matrix converge.

The Gaussian lift proposition gives the
sectorwise limit. The overlap formula
gives its Hilbert limit after applying \(\pi\).

The substitutions on \(y,C\) give
\eqref{v4-eq:singular-gaussian-pair-law}.
The resolvent identity gives the same
composition for \(K'\).
For the trace update, subtract the
purely linear transform of \(K\) from
the final response.
Insert the intermediate response between
them and take the partial trace.
The two trace-class differences telescope,
and finite mode matrices pass outside
the partial trace. This proves the trace law.

It remains to check the scalar, including
the sign of its regularized determinant.
Put
\[
 A=\widetilde B_2K,\qquad
 D=\widetilde T_2\widetilde B_1\widetilde T_2^t,\qquad Q=DK.
\]
By similarity, the second determinant
has argument \(1-Q(1-A)^{-1}\).
The exact factorization is
\[
                  (1-Q(1-A)^{-1})(1-A)=1-A-Q.
\]

For Hilbert--Schmidt \(U,V\), the regularized
product identity is
\[
 \det{}_2((1-U)(1-V))
       =\det{}_2(1-U)\det{}_2(1-V)e^{-\operatorname{Tr}(UV)}.
\]
In finite rank this follows by separating
the linear trace from each ordinary determinant.
Approximation in Hilbert--Schmidt norm extends
it, since \(UV\) converges in trace norm.

This is also the \(p=2\) product formula
for regularized Fredholm determinants.\footnote{
T.~Britz, A.~Carey, F.~Gesztesy, R.~Nichols,
F.~Sukochev and D.~Zanin.
\emph{The product formula for regularized Fredholm determinants},
equations~(1.3)--(1.5),
\href{https://arxiv.org/abs/2007.12834}{arXiv:2007.12834}.}

The two inverse square roots therefore
leave the multiplier
\[
                \exp\left[-\frac12
                     \operatorname{Tr}\bigl(Q(1-A)^{-1}A\bigr)\right].
\]
The trace update contributes its inverse:
\[
 \begin{aligned}
 \frac12\operatorname{tr}_d(B_1\tau_2)
   -\frac12\operatorname{tr}_d(T_2B_1T_2^t\tau)
     &=\frac12\operatorname{Tr}\bigl(D(K^{B_2}-K)\bigr)\\
     &=\frac12\operatorname{Tr}\bigl(Q(1-A)^{-1}A\bigr).
 \end{aligned}
\]
All traced products are trace class.
The exponents cancel, proving
\eqref{v4-eq:singular-gaussian-scalar-law}.
Associativity of the matrix pair law
then proves every finite regrouping on
its common domain.
\end{proof}

For \(d=1\), \(T=1\) and \(B=t\),
the construction becomes
\[
 K_t=K(1-tK)^{-1},\qquad
 \tau_t=\tau+\operatorname{Tr}(K_t-K),\qquad
 F_t=e^{t\tau/2}\det{}_2(1-tK)^{-1/2}.
\]
Differentiating the convergent series gives
\[
       \dot K_t=K_t^2,\qquad
       \dot\tau_t=\operatorname{Tr}(K_t^2),\qquad
       \frac{d}{dt}\log F_t=\tau_t/2.
\]
The trace-class difference replaces an
undefined \(\operatorname{Tr}K\).
Forgetting \(\tau\) recovers the defect
in \eqref{v4-eq:singular-determinant-defect}.

Two trace data \(\tau,\widehat\tau\) for
the same \(K\) project to the same
Hilbert vector \(S_K\).
Their transformed Hilbert limits instead
differ by
\[
               \exp\left[\frac12
                     \operatorname{tr}_d(B(\tau-\widehat\tau))\right].
\]
If this scalar is not one, no continuous
operator on the old Hilbert space can
realize both cutoff limits.
Signed reservoirs do not alter the divergence
of the original positive-cutoff sequence
for a positive non-trace-class kernel.
They specify a different trace datum.

Even the response endpoints do not
determine the scalar. On one occupied
oscillator take \(0<\kappa<1\),
\(K=\tau=\kappa\), \(T=a>0\) near one,
and \(B=(1-a^2)/\kappa\). Then
\[
                      K'=\kappa,\qquad F=a^{-1}.
\]
The coordinate germ
\(f(z)=az/\sqrt{1-2Bz^2}\) has these
first-mode data. The coordinate arrow
and its normalization must therefore
be retained even when the response returns
to its initial kernel.

These formulas define a Gaussian domain
in \(\prod_E\mathcal W_E\).
They do not extend every coordinate change
to that unrestricted product.
Coordinate contractions can lower energy
without bound. The Gaussian resolvent is
what supplies convergence on this domain.

\subsection{Unitary implementation is a separate condition}

Let a bounded invertible real symplectic
map act on a one-particle Hilbert space.
Relative to its chosen complex structure,
write \(\alpha\) for its linear block
and \(\beta\) for its antilinear block.
Shale's criterion says that its action
on the canonical commutation relations
has a unitary Fock implementer exactly
when \(\beta\) is Hilbert--Schmidt.\footnote{
L.~Bruneau and J.~Dereziński,
\emph{Bogoliubov Hamiltonians and one-parameter groups
of Bogoliubov transformations}.
Definition~3.1, Assumption~3.A and Theorem~3.2,
\href{https://arxiv.org/abs/math-ph/0511069}{arXiv:math-ph/0511069}.}
A formal coordinate germ need not define
such a bounded symplectic Hilbert map.

If the mode action is identical in every
species, then
\(\beta=\beta_{\mathrm{mode}}\otimes1_H\).
For a unit mode vector \(v\) with
\(\beta_{\mathrm{mode}}v\ne0\),
\[
 \sum_i\|\beta(v\otimes e_i)\|^2
            =\sum_i\|\beta_{\mathrm{mode}}v\|^2=\infty.
\]
Thus this action has no unitary implementer.
When \(\beta=0\) and \(\alpha\) is unitary,
ordinary second quantization does provide one.

For the uniform one-mode squeeze
\(\alpha=5\,1_H/4\), \(\beta=3\,1_H/4\),
the finite-species kernel is \(K_N=3P_N/5\), and
\[
 \|\beta_N\|_2^2=9N/16,\qquad
 \langle\mathbf1,\Omega_{K_N}\rangle=(16/25)^{N/4}\longrightarrow0.
\]
This distinguishes unitary implementation
from convergence of selected preparations
and from covariance of finite-energy fields.

A continuous prepared functional on a
bounded-energy projective sector can be
transported by precomposition with \(R_f\).
Preservation of positivity additionally
requires its adjoint and positive cone
to be preserved. This covector construction
alone supplies neither a global state
nor a physical sewing map.

\section{State families and completed field operations}
\label{v4-sec:singular-state-families}

The preceding completions control species tensors
or specified Gaussian families. A formal parameter
raises a different question. When may that
parameter be identified with an insertion coordinate?
An arbitrary formal series does not suffice.

Return to a characteristic-zero field \(k\)
with discrete coefficients. For the current
\(J\in H_\kappa\), introduce a parameter
\(\xi\), independent of the insertion coordinate \(u\).
The state family
\[
 A(\xi)=\sum_{r\geq0}
             \frac{(-1)^r\xi^r}{(r+1)!}T^rJ
\]
exists in \(H_\kappa[[\xi]]\).
Its coefficientwise field on \(J\) has vacuum
polar part
\[
                         \kappa\sum_{r\geq0}\xi^ru^{-r-2}\mathbf1.
\]

The partial state sums converge \(\xi\)-adically.
After the substitution \(\xi=u\), their
\(u^{-2}\mathbf1\) coefficient is \(N\kappa\)
for a sum of \(N\) terms. It does not stabilize.
Thus this substitution is not continuous
on all formal state families with discrete scalars.

\subsection{The exact domain for diagonal evaluation}

Let \(V\) be a vertex algebra over \(k\).
The same discussion permits a differential
graded vertex algebra, with all series taken
in each cohomological degree.
Give \(V((u))\) the topology with neighborhoods
\(u^qV[[u]]\), \(q\in\mathbb Z\).
For \(a,b\in V\), put
\[
 \nu_b(a)=\operatorname{ord}_uY(a,u)b,\qquad
                         \operatorname{ord}_u0=+\infty.
\]

Define
\begin{equation}\label{v4-eq:singular-admissible-family}
 \mathscr A_Y=
   \left\{\sum_{r\geq0}\xi^ra_r:
          r+\nu_b(a_r)\longrightarrow+\infty
                             \text{ for every }b\in V\right\}.
\end{equation}
The parameter \(\xi\) has cohomological degree zero.
For a finite set \(B\subset V\) and \(q\in\mathbb Z\),
let \(U_{B,q}\) require
\[
                   r+\nu_b(a_r)\geq q
                   \quad\text{for every }r\geq0,\ b\in B.
\]
These subspaces define a linear topology on
\(\mathscr A_Y\).
Equip \(\operatorname{Hom}_k(V,V((u)))\)
with pointwise Laurent convergence.

\begin{proposition}\label{v4-prop:singular-diagonal-evaluation}
Polynomial families are dense in \(\mathscr A_Y\).
Their diagonal evaluation extends uniquely to
a continuous map
\[
 \operatorname{ev}_Y:\mathscr A_Y
                 \longrightarrow\operatorname{Hom}_k(V,V((u))),
 \qquad
 \operatorname{ev}_Y(A)(u)b=\sum_{r\geq0}u^rY(a_r,u)b.
\]
The limit condition in
\eqref{v4-eq:singular-admissible-family}
is necessary and sufficient for convergence
on every input.
The domain is preserved by coefficient
translation \(T\), by \(\partial_\xi\),
and by the differential in the graded case.
Evaluation commutes with that differential and satisfies
\[
 \operatorname{ev}_Y((T+\partial_\xi)A)
                         =\partial_u\operatorname{ev}_Y(A),
 \qquad
 [T,\operatorname{ev}_Y(A)]=\operatorname{ev}_Y(TA).
\]
\end{proposition}

\begin{proof}
The limit condition makes the series
converge in the complete Laurent space
on each fixed input. Conversely, successive
differences of convergent partial sums
must tend to zero, giving the same condition.
The image of \(U_{B,q}\) sends every
\(b\in B\) into \(u^qV[[u]]\).
This proves continuity.

Truncation tails tend to zero in every
\(U_{B,q}\), so polynomials are dense.
The target is Hausdorff, proving uniqueness.
The source is Hausdorff as well:
the vacuum identity gives \(\nu_{\mathbf1}(a)=0\)
for every nonzero \(a\).

Translation covariance gives
\(\nu_b(Ta)\geq\nu_b(a)-1\).
Reindexing gives the same loss of at most
one for \(\partial_\xi\).
For homogeneous \(a\), the differential identity gives
\[
                 \nu_b(da)\geq\min\{\nu_b(a),\nu_{db}(a)\}.
\]
These estimates preserve the domain and
justify termwise differentiation.

The product rule on \(u^rY(a_r,u)\)
gives the first translation identity.
Summing \([T,Y(a_r,u)]=Y(Ta_r,u)\)
gives the second.
The graded Hom differential gives the
differential assertion.
\end{proof}

There are admissible infinite families
with unbounded derivative order.
For example,
\[
 \widetilde A(\xi)
      =(1+\xi)\sum_{r\geq0}
                \frac{(-1)^r\xi^{2r}}{(r+1)!}T^rJ
\]
belongs to \(\mathscr A_Y\), since
\(\nu_b(T^rJ)\geq\nu_b(J)-r\).
Its vacuum coefficient on \(J\) after evaluation is
\[
               \kappa u^{-2}(1+u)\sum_{r\geq0}u^r
                         =\kappa u^{-2}\frac{1+u}{1-u}.
\]
This is expanded at \(u=0\).
The parameter degree now outruns the
derivative pole order.

\subsection{A weight completion closed under every mode}

Suppose \(V\) also has a nonnegative
integer weight grading with
\[
 V=\bigoplus_{h\geq0}V_h,\quad
 \mathbf1\in V_0,\quad T(V_h)\subset V_{h+1},\quad
 d(V_h)\subset V_h,\quad
 (V_h)_{(n)}V_{h'}\subset V_{h+h'-n-1}.
\]
Negative weight spaces are zero.
Weight is independent of cohomological degree,
and its spaces need not be finite dimensional.
Let \(w(a)\) be the largest weight of a
nonzero component of \(a\), with \(w(0)=-\infty\).

Put
\begin{equation}\label{v4-eq:singular-weight-completion}
 \mathscr A_{\mathrm{wt}}(V)
   =\left\{\sum_{r\geq0}\xi^ra_r:
                r-w(a_r)\longrightarrow+\infty\right\}.
\end{equation}
Define
\[
 \gamma(A)=\inf_{r\geq0}\{r-w(a_r)\},\qquad
 \mathcal F^q=\{A:\gamma(A)\geq q\},\qquad
 \gamma(0)=+\infty.
\]
The \(\mathcal F^q\) are the neighborhoods
of zero. Each sum is still taken in
a fixed cohomological degree.

On polynomials \(W=V\otimes k[\xi]\),
use the tensor-product fields
\begin{equation}\label{v4-eq:singular-parameter-field}
 Y_W(\xi^ra,z)(\xi^sb)
                    =(\xi+z)^r\xi^sY_V(a,z)b.
\end{equation}
The second factor has its commutative
differential-algebra vertex structure.
Translation on \(W\) is \(D=T+\partial_\xi\).
In an iterated field, the shifted
parameter becomes \(\xi+z_2+z_0\).
This equals \(\xi+z_1\) under
\(z_1=z_2+z_0\). Thus the tensor product
retains the vertex identities.

\begin{theorem}\label{v4-thm:singular-completed-modes}
The space \(\mathscr A_{\mathrm{wt}}(V)\)
is the Hausdorff completion of \(W\)
for the degree \(\deg(\xi^ra)=r-h\),
\(a\in V_h\).
Every mode extends to the continuous
bilinear product
\begin{equation}\label{v4-eq:singular-completed-mode}
 A_{[n]}B
    =\sum_{r,s\geq0}\sum_{j=0}^r
       \binom rj\xi^{r+s-j}(a_r)_{(n+j)}b_s.
\end{equation}
Each parameter coefficient is a finite sum,
and
\begin{equation}\label{v4-eq:singular-completed-mode-bound}
 \gamma(A_{[n]}B)\geq\gamma(A)+\gamma(B)+n+1,
 \qquad
 \gamma(A_{[n]}B)-n\longrightarrow+\infty
                         \quad(n\to+\infty).
\end{equation}
The vacuum, differential and translation
extend with their usual identities.
For homogeneous \(A,B,C\) and integers \(p,q,n\),
the three series in the identity
\begin{align}
 \sum_{j\geq0}\binom pj(A_{[n+j]}B)_{[p+q-j]}C
     &=
       \sum_{j\geq0}(-1)^j\binom nj
                       A_{[p+n-j]}(B_{[q+j]}C)
                  \label{v4-eq:singular-completed-Borcherds}\\
     &\quad-
       (-1)^{n+|A||B|}
       \sum_{j\geq0}(-1)^j\binom nj
                       B_{[q+n-j]}(A_{[p+j]}C)
                  \nonumber
\end{align}
converge and satisfy the displayed equality.
Every fixed coefficient of \(\xi\) receives
only finitely many summands from each series.
\end{theorem}

\begin{proof}
Decompose the polynomial space as
\[
 W=\bigoplus_{\delta\in\mathbb Z}W_\delta,\qquad
 W_\delta=\bigoplus_{\substack{r,h\geq0\\r-h=\delta}}\xi^rV_h.
\]
Its completion consists of sums bounded
below in \(\delta\), with finitely many
pairs \((r,h)\) below every degree cutoff.
This is equivalent to the limit condition
in \eqref{v4-eq:singular-weight-completion}.
For fixed \(r\), every weight satisfies
\(r-h\leq r\), so its coefficient remains
a finite sum in \(V\).
This proves completeness, Hausdorffness
and density of the polynomial subspace.

Expanding \eqref{v4-eq:singular-parameter-field}
gives the mode formula on polynomials.
A nonzero summand at parameter exponent
\(\ell=r+s-j\) has weight at most
\(w(a_r)+w(b_s)-n-j-1\). Hence
\begin{equation}\label{v4-eq:singular-mode-surplus}
 \ell-w((a_r)_{(n+j)}b_s)
       \geq(r-w(a_r))+(s-w(b_s))+n+1.
\end{equation}
Nonnegative output weight also gives
\[
              (r-w(a_r))+(s-w(b_s))\leq\ell-n-1.
\]

The two sequences on the left tend
to infinity and are bounded below.
Only finitely many pairs \((r,s)\)
can contribute to a fixed \(\ell\).
Then \(j\) is determined.

As \(\ell\to\infty\), any contributing
pair satisfies \(r+s\geq\ell\).
Its sum of surpluses therefore tends
to infinity. This proves that the
product belongs to the completion,
with the first bound in
\eqref{v4-eq:singular-completed-mode-bound}.
The bound also proves continuity.

Approximate \(A,B\) by polynomials \(A_0,B_0\)
with errors of arbitrarily large \(\gamma\).
For all sufficiently large \(n\),
\((A_0)_{[n]}B_0=0\).
Apply the first bound to
\[
           (A-A_0)_{[n]}B+(A_0)_{[n]}(B-B_0).
\]
After subtracting \(n\), its lower
bound can be made arbitrarily large,
uniformly in \(n\). This proves the
second limit.

The differential preserves
the filtration. \(D\) lowers it by at most one.
Thus their polynomial identities extend continuously.

Each summand in the three iterated
products of \eqref{v4-eq:singular-completed-Borcherds}
has filtration lower bound
\[
                       \gamma(A)+\gamma(B)+\gamma(C)+p+q+n+2.
\]
This bound is independent of \(j\).
Replacing any input by a high-filtration
error raises it uniformly.
Approximate all three inputs by polynomials
with bounded-below filtrations.
Their three series are finite, by
ordinary mode truncation.
Beyond the last polynomial term, every
completed summand therefore lies in any
prescribed \(\mathcal F^N\), once the
approximations are sufficiently close.

The series converge and inherit the
polynomial Borcherds identity.

Finally, an element of \(\mathcal F^{L+1}\)
has zero \(\xi^L\) coefficient, since
that coefficient would have negative weight.
The preceding convergence proof places
all sufficiently late summands there.
This proves the asserted finiteness
of each parameter coefficient.
\end{proof}

\subsection{Modules and the lower weight bound}

Let \(M=\bigoplus_{\ell\geq b_0}M_\ell\)
be a weight-graded \(V\)-module, where
\(b_0\in\mathbb Z\), the differential
preserves weight, and
\[
                       (V_h)_{(m)}M_\ell\subset M_{h+\ell-m-1}.
\]
Equip \(\operatorname{Hom}_k(M,M((u)))\)
with pointwise Laurent convergence.

\begin{proposition}\label{v4-prop:singular-completed-module}
Polynomial evaluation extends continuously to
\[
 \widehat Y_M:\mathscr A_{\mathrm{wt}}(V)
                \longrightarrow\operatorname{Hom}_k(M,M((u))),
 \qquad
 \widehat Y_M(A,u)b=\sum_{r\geq0}u^rY_M(a_r,u)b.
\]
Its mode on \(b\in M_\ell\) is the finite sum
\begin{equation}\label{v4-eq:singular-completed-module-mode}
                       A^M_{[m]}b
                              =\sum_{r\geq0}(a_r)_{(m+r)}b,
\end{equation}
and
\[
                  \operatorname{ord}_u\widehat Y_M(A,u)b
                                      \geq\gamma(A)+b_0-\ell.
\]
The module identities hold with the completed
products \eqref{v4-eq:singular-completed-mode}.
Their sums are locally finite on each fixed
module vector. The action is natural in
strict weight-preserving module maps,
commutes with the differential, and
intertwines \(D\) with \(\partial_u\).
\end{proposition}

\begin{proof}
The \(r\)-th summand of
\eqref{v4-eq:singular-completed-module-mode}
has weight at most
\[
                         \ell-m-1-(r-w(a_r)).
\]
It vanishes when this is less than \(b_0\).
The surplus limit proves finiteness and
the common Laurent lower bound.
That same bound proves continuity.
On polynomials this is the tensor-product
module where \(f(\xi)\) acts by \(f(u)\).

In a module Borcherds identity, each
composite summand sends a weight-\(\ell\)
input to weights at most
\[
                       \ell-p-q-n-2-\gamma(A)-\gamma(B),
\]
independently of the summation index.
A sufficiently high-filtration error in
either input annihilates every summand.
Replace both inputs by polynomials at
that precision. The polynomial identity
is finite on the given vector and
agrees term by term with the completed one.
This proves local finiteness and the identity.
The differential, naturality and translation
statements follow from the polynomial
identities by the same estimate.
\end{proof}

The regular module gives
\(\mathscr A_{\mathrm{wt}}(V)\subset\mathscr A_Y\).
It also contains the derivative-controlled
families with
\[
 a_r\in\sum_{j=0}^{d_r}T^jE,\qquad
                         r-d_r\longrightarrow+\infty,
\]
where \(E\) is a fixed finite-dimensional
state space. If \(c\) is its maximum
weight, then \(w(a_r)\leq c+d_r\).
The stated inclusion follows directly.

The lower module bound is substantial.
Consider the weak Heisenberg module
\[
 M=k[y,x_2,x_3,\ldots],\qquad
                         w(y)=-1,\quad w(x_m)=m,
\]
with operators
\begin{equation}\label{v4-eq:singular-unbounded-module}
 J_1=y,\qquad J_{-1}=-\kappa\partial_y,\qquad J_0=0,
 \qquad
 J_m=m\kappa\partial_{x_m},\quad J_{-m}=x_m\quad(m\geq2).
\end{equation}
The identities
\([y,-\kappa\partial_y]=\kappa\) and
\([m\kappa\partial_{x_m},x_m]=m\kappa\)
give the current commutators.
Each polynomial is annihilated by sufficiently
large positive modes. The normally ordered
derivative fields therefore define a weak
\(H_\kappa\)-module.

At each output power only finitely many
negative-mode indices can contribute,
after the finitely many possible positive
modes have acted. Wick expansion proves
the module field identities on this
polynomial space.

Put \(a_r=J_{-1}^r\mathbf1\) in the
state algebra and
\[
                  A(\xi)=\sum_{r\geq0}\xi^{2r}a_r,
 \qquad
                  A_N(\xi)=\sum_{r=0}^N\xi^{2r}a_r.
\]
The state weight is \(w(a_r)=r\).
Thus \(A_N\to A\) in the weight completion,
with \(\gamma(A-A_N)=N+1\).
On the constant polynomial \(1\in M\),
\begin{equation}\label{v4-eq:singular-unbounded-evaluation}
       [u^0]\sum_{r=0}^N u^{2r}Y_M(a_r,u)1
                                      =\sum_{r=0}^Ny^r.
\end{equation}

Indeed \(Y_M(a_r,u)=:J(u)^r:\).
After normal ordering, all positive modes
except \(J_1=y\) annihilate \(1\).
The coefficient of \(u^{-2r}\) comes only
from \(r\) copies of \(yu^{-2}\).
Every negative-mode factor has a nonnegative
power of \(u\) and gives a larger exponent.
This proves the displayed formula.

Its successive constant coefficients differ
by \(y^{N+1}\ne0\). The evaluations are
not Cauchy modulo \(uM[[u]]\).
Consequently no continuous extension of
polynomial evaluation exists on this
module with the discrete Laurent target.
The constant polynomial \(1\in M\)
is not a highest-weight vacuum:
\(J_1 1=y\).
The obstruction concerns this module
and target, while the lower-bounded
action remains defined.

\subsection{Common-scale fields and the geometric boundary}

For \(n\) insertions put \(z_i=tx_i\)
with \(x_1=1\), and let
\[
 R_n=k[x_2,\ldots,x_n]
          [\,x_i^{-1},(x_i-x_j)^{-1}:
                    2\leq i\leq n,\ 1\leq i<j\leq n\,].
\]
The first inverses keep the insertions
away from zero. The second keep their
positions distinct.

For homogeneous states of weights \(h_i\)
and an input \(b\in M_\ell\), the
simultaneous rational product has its
\(t^d\) coefficient in
\[
                         M_{\ell+\sum_i h_i+d}\otimes R_n.
\]
To obtain it, clear the finitely many
pairwise locality poles and use the
module field identities to compare orders.
Place each insertion in turn at the
rightmost position. Lower truncation
on \(b\) then gives a lower bound for
that variable's exponent.
The pole-clearing factors are polynomials,
so a common lower bound survives in
every variable.

The resulting regular coefficients have
fixed total degree and finitely many
monomials at each degree.
Dividing by the cleared factors gives
the same rational coefficient in every
chamber. The weight bound makes its
\(t\)-order at least \(b_0-\ell-\sum_i h_i\).

For families \(A_i=\sum_{r\geq0}\xi^ra_{i,r}\),
their term of parameter degrees \(r_i\)
therefore has order at least
\[
                 b_0-\ell+\sum_i(r_i-w(a_{i,r_i})).
\]
Only finitely many tuples lie below
any cutoff. Hence these simultaneous
fields converge in
\[
                                  (M\otimes R_n)((t))
\]
for all \(A_i\in\mathscr A_{\mathrm{wt}}(V)\).
This gives a common-scale evaluation
on every lower-bounded module.

The completed products need not have
finite mode truncation. In \(H_\kappa\),
take
\[
                         A(\xi)=\sum_{r\geq0}\frac{\xi^{2r}}{r!}T^rJ.
\]
Its coefficient of \(\xi^{2r}\)
has weight \(r+1\), so it belongs
to the weight completion.
For \(n\geq1\),
\begin{equation}\label{v4-eq:singular-infinite-current-modes}
 A_{[n]}J
    =\kappa\sum_{r\geq n-1}(-1)^r(r+1)
                  \binom{2r}{r+1-n}\xi^{r+n-1}\mathbf1.
\end{equation}

This follows from
\((T^rJ)_{(r+1)}J=(-1)^r(r+1)!\kappa\mathbf1\)
and the completed mode formula.
Its first term is
\((-1)^{n-1}n\kappa\xi^{2n-2}\mathbf1\).
Thus every displayed mode is nonzero, while
\[
                         \gamma(A_{[n]}J)-n=n-2\longrightarrow+\infty.
\]
The mode identities converge despite
the absence of finite singular-mode truncation.

The same example exhibits a displaced
collision in its common-scale field.
Its vacuum coefficient with \(J(z_2)\) is
\[
 \sum_{r\geq0}\frac{z_1^{2r}}{r!}
       \partial_{z_1}^r\frac{\kappa}{(z_1-z_2)^2}
                    =\frac{\kappa}{(z_1+z_1^2-z_2)^2}.
\]
For \(z_1=t\), \(z_2=tx_2\), the right
side is expanded as
\(\kappa t^{-2}(1-x_2+t)^{-2}\) in \(R_2((t))\).
Its displaced pole explains why closed
completed modes do not automatically define
residue maps on the original collision carrier.

\begin{question}\label{v4-q:singular-completed-geometric-comparison}
For a specified completion and its module
domain, construct collision extensions and
normal residue maps compatible with the
completed field products and differential-module
coefficients. For independently defined arithmetic
coefficients, construct a supported comparison
preserving all collision channels, both vacuum
maps and the prepared functional.
Which geometric data determine compatible
bidifferentials and trace matrices, and
which comparisons preserve the scalar
under the prescribed sewing of states?
\end{question}

The algebraic Heisenberg background already
has its cut bar comparison and framed
field recovery. Its analytic completions
retain the operations proved on their
specified domains above.
An arithmetic use must additionally construct
the coefficient and state comparison in
the question. Positivity of an independently
defined Weil form, a derived factorizing
center and Hilbert unitary implementation
remain separate requirements.
